\documentclass[11pt]{article}
\usepackage[margin=1in]{geometry}
\usepackage{amsmath, amssymb, amsthm}
\usepackage{cancel}
\usepackage{enumitem}
\usepackage{xcolor}
\usepackage{hyperref}
\hypersetup{colorlinks=true, linkcolor=blue!50!black}

\theoremstyle{definition}
\newtheorem{defn}{Definition}[section]
\newtheorem{thm}[defn]{Theorem}
\newtheorem{ex}[defn]{Example}
\newtheorem{method}[defn]{Method}
\newtheorem{lem}[defn]{Lemma}
\newtheorem{prob}{Practice Problem}
\newtheorem{mprob}{Exam Problem}

\newcommand{\N}{\mathbb{N}}
\newcommand{\Z}{\mathbb{Z}}
\newcommand{\Q}{\mathbb{Q}}
\newcommand{\R}{\mathbb{R}}
\newcommand{\eps}{\varepsilon}
\newcommand{\key}[1]{\textcolor{red!70!black}{\textbf{#1}}}

\newenvironment{scratch}{\par\medskip\noindent\textit{Scratch work (not part of the proof):}\ }{\par\medskip}

\title{MATH 3150 --- Exam 1 Study Guide\\
\large Methods, templates, worked examples, and practice problems\\
\normalsize (Weeks 1--5: $\N,\Q,\R$, induction, sup/inf, completeness,
absolute value, limits of sequences, monotone sequences, subsequences,
Cauchy sequences)}
\author{}
\date{Exam 1: Tuesday, October 13, in class (12:30--1:45, MONT 225)}

\begin{document}
\maketitle
\tableofcontents
\newpage

% =====================================================================
\section{How to use this guide}

Each section follows the same pattern: \textbf{definitions you must be able
to state word for word}, then \textbf{the method}, then \textbf{worked
examples}. Section~\ref{sec:practice} has practice problems, and their
solutions are in Section~\ref{sec:solutions}. Try each problem on paper
before you look at the solution.

\medskip
\noindent\key{Exam 1 scope.} The syllabus says Weeks 1--4, as tested in
HW1--HW4. \textbf{Cauchy sequences will also be tested}, so
Section~\ref{sec:cauchy} covers them in full, along with the subsequence
facts that the Cauchy proofs need.
\begin{itemize}[nosep]
  \item HW1: set identities, induction, rational/irrational arguments.
  \item HW2: density of $\Q$ and of the irrationals, sup/inf, the
        completeness axiom, $\sup(A+B)$.
  \item HW3: triangle and reverse triangle inequality, $\eps$--$N$ limits,
        $x_n\to\infty$, monotone sequences.
  \item HW4 and Week 4 lectures: limit laws, basic limits, limits equal to
        $\pm\infty$, the Monotone Convergence Theorem.
  \item Quiz 1: $\sup\{xy\}=\sup A\cdot\sup B$ for positive sets
        (worked in full in Example~\ref{ex:quiz1}).
  \item Week 5: subsequences, Bolzano--Weierstrass, \textbf{Cauchy
        sequences}, and the \textbf{Nested Interval Theorem}
        (Section~\ref{sec:nit}).
  \item Not tested: the Week 6 topology material.
\end{itemize}

\medskip
\noindent\key{Exam vs.\ quiz.} The quiz was one proof in 20 minutes. The
exam is 75 minutes and will probably have several problems: expect one
or two definitions or statements, a sup/inf proof, an $\eps$--$N$ proof,
a monotone or recursive sequence, and a Cauchy question.
Section~\ref{sec:mock} is a practice exam in that format.

\medskip
\noindent\key{The single most important skill:} match the \emph{shape}
of the claim to the proof technique.

\begin{center}
\renewcommand{\arraystretch}{1.3}
\begin{tabular}{|p{6cm}|p{8.5cm}|}
\hline
\textbf{Claim looks like\ldots} & \textbf{Reach for\ldots}\\ \hline
$S=T$ (sets) & Double inclusion, element chasing \\ \hline
``for all $n\ge n_0$'' & Induction \\ \hline
``$x$ is irrational'', ``no such $\ldots$'' & Contradiction \\ \hline
``is it always true?'' (suspect false) & One counterexample \\ \hline
$M=\sup S$ & (i) $M$ is an upper bound, (ii) $\forall\eps>0\ \exists s\in S$, $s>M-\eps$ \\ \hline
$\sup A \le$ something & Show that ``something'' is an upper bound of $A$ \\ \hline
$\lim x_n = L$ & $\eps$--$N$: scratch work, then a forward write-up \\ \hline
$\lim x_n=\infty$ & $\forall M\ \exists N$: $x_n>M$ for $n\ge N$ \\ \hline
$x_n$ converges (limit unknown) & Monotone Convergence Theorem \\ \hline
$|a|\le c$ & Show $-c\le a\le c$ \\ \hline
\end{tabular}
\end{center}

\subsection*{Definitions cheat sheet (memorize word for word)}
\begin{center}
\renewcommand{\arraystretch}{1.35}
\begin{tabular}{|p{3.4cm}|p{11.2cm}|}
\hline
Upper bound & $M$ with $s\le M$ for all $s\in S$\\ \hline
$\sup S$ & upper bound $M$ such that no $M'<M$ is an upper bound
$\iff$ upper bound and $\forall\eps>0\ \exists s\in S:\ s>M-\eps$\\ \hline
$\inf S$ & lower bound $m$ such that no $m'>m$ is a lower bound $\iff$
lower bound and $\forall\eps>0\ \exists s\in S:\ s<m+\eps$\\ \hline
Completeness & nonempty $+$ bounded above $\Rightarrow$ $\sup$ exists in $\R$\\ \hline
Archimedean & $\forall x\in\R\ \exists n\in\N:\ n>x$; equivalently $\forall\eps>0\ \exists n:\ \frac1n<\eps$\\ \hline
Density & $x<y\Rightarrow\exists r\in\Q$ and $\exists z\notin\Q$ with $x<r,z<y$\\ \hline
$|x|\le a$ & $\iff -a\le x\le a$ (for $a\ge0$)\\ \hline
$x_n\to L$ & $\forall\eps>0\ \exists N\ \forall n\ge N:\ |x_n-L|<\eps$\\ \hline
$x_n\not\to L$ & $\exists\eps>0\ \forall N\ \exists n\ge N:\ |x_n-L|\ge\eps$\\ \hline
$x_n\to\infty$ & $\forall M\ \exists N\ \forall n\ge N:\ x_n>M$\\ \hline
Bounded seq. & $\exists M>0\ \forall n:\ |x_n|\le M$\\ \hline
Increasing & $x_n\le x_{n+1}$ for all $n$ \quad (decreasing: $x_n\ge x_{n+1}$)\\ \hline
MCT & bounded $+$ monotone $\Rightarrow$ convergent (to $\sup$ or $\inf$ of the terms)\\ \hline
Cauchy & $\forall\eps>0\ \exists N\ \forall n,m\ge N:\ |x_n-x_m|<\eps$\\ \hline
Not Cauchy & $\exists\eps>0\ \forall N\ \exists n,m\ge N:\ |x_n-x_m|\ge\eps$\\ \hline
Subsequence & $(x_{n_k})_k$ with $n_1<n_2<n_3<\cdots$ (so $n_k\ge k$)\\ \hline
Cauchy criterion & in $\R$: $(x_n)$ converges $\iff$ $(x_n)$ is Cauchy\\ \hline
B--W & every bounded sequence has a convergent subsequence\\ \hline
NIT & closed, bounded, nested $[a_n,b_n]$ $\Rightarrow$ $\bigcap_n[a_n,b_n]\ne\emptyset$; a single point if $b_n-a_n\to0$\\ \hline
\end{tabular}
\end{center}

% =====================================================================
\section{Logic and proof techniques}

\subsection{The basic proof moves}
\begin{itemize}[nosep]
  \item \textbf{Direct:} assume the hypotheses and chain implications to
        the conclusion.
  \item \textbf{Contrapositive:} $P\Rightarrow Q$ is equivalent to
        $\neg Q\Rightarrow\neg P$. Example: ``$n^2$ even $\Rightarrow$ $n$ even''
        is easiest as ``$n$ odd $\Rightarrow$ $n^2$ odd''.
  \item \textbf{Contradiction:} assume $\neg(\text{claim})$ and derive
        something false.
  \item \textbf{Proving ``$a\le b$'' via $\eps$:} if $a<b+\eps$ for every
        $\eps>0$, then $a\le b$. (Otherwise take $\eps=a-b>0$.) Likewise,
        if $|a-b|<\eps$ for every $\eps>0$, then $a=b$. \key{This is how
        uniqueness of limits and many sup facts are proved.}
  \item \textbf{``For every'' statements:} begin with ``Let \ldots\ be
        arbitrary.'' \textbf{``There exists'' statements:} construct or
        exhibit the object.
\end{itemize}

\subsection{Negating quantifiers}
To negate a statement, swap $\forall\leftrightarrow\exists$ and negate the
inner condition.
\[
\neg\bigl(\forall \eps>0\ \exists N\ \forall n\ge N:\ |x_n-L|<\eps\bigr)
\iff
\exists \eps>0\ \forall N\ \exists n\ge N:\ |x_n-L|\ge\eps .
\]
So ``$x_n\not\to L$'' means: \emph{some} $\eps$ is a bad tolerance, and
the sequence keeps returning to distance $\ge\eps$ from $L$ no matter how
far out you go. You use this to prove divergence (Example~\ref{ex:diverge}).

\subsection{Double inclusion (HW1 P1)}
\begin{method}
To prove $S=T$, prove $S\subseteq T$ and $T\subseteq S$. For
$S\subseteq T$: ``Let $x\in S$ be arbitrary. \ldots\ Therefore $x\in T$.''
Split into cases wherever the definition involves ``or'' (a union) or
``not (and)'' (a complement of an intersection).
\end{method}

\begin{ex}[De Morgan-type identity]
$A\setminus(B\cup C)=(A\setminus B)\cap(A\setminus C)$.
\begin{proof}[Quiz-length version]
($\subseteq$) Let $x\in A\setminus(B\cup C)$. Then $x\in A$ and
$x\notin B\cup C$, so $x\notin B$ \emph{and} $x\notin C$. Hence
$x\in A\setminus B$ and $x\in A\setminus C$.

($\supseteq$) Let $x\in(A\setminus B)\cap(A\setminus C)$. Then $x\in A$,
$x\notin B$, and $x\notin C$, so $x\notin B\cup C$. Hence
$x\in A\setminus(B\cup C)$.
\end{proof}

\begin{proof}[Full proof, step by step]
\emph{Step 1 (unpack the definitions you will use).}
\[
x\in A\setminus D\iff (x\in A\text{ and }x\notin D),\qquad
x\in D\cup E\iff(x\in D\text{ or }x\in E).
\]
Negating the second one gives
$x\notin D\cup E\iff(x\notin D\text{ and }x\notin E)$. This negation of an
``or'' is the whole content of the identity.

\emph{Step 2 ($\subseteq$: start with an arbitrary element).} Let
$x\in A\setminus(B\cup C)$ be arbitrary. By the definition of set
difference, $x\in A$ and $x\notin B\cup C$.

\emph{Step 3 (negate the union).} By Step 1, $x\notin B\cup C$ means
$x\notin B$ and $x\notin C$.

\emph{Step 4 (repackage).} From $x\in A$ and $x\notin B$ we get
$x\in A\setminus B$. From $x\in A$ and $x\notin C$ we get
$x\in A\setminus C$. So $x$ is in both, i.e.\
$x\in(A\setminus B)\cap(A\setminus C)$. Since $x$ was arbitrary,
$A\setminus(B\cup C)\subseteq(A\setminus B)\cap(A\setminus C)$.

\emph{Step 5 ($\supseteq$: arbitrary element of the right side).} Let
$x\in(A\setminus B)\cap(A\setminus C)$ be arbitrary. By the definition of
intersection, $x\in A\setminus B$ and $x\in A\setminus C$, so $x\in A$,
$x\notin B$, and $x\notin C$.

\emph{Step 6 (rebuild the union).} Since $x\notin B$ and $x\notin C$,
Step 1 gives $x\notin B\cup C$. Together with $x\in A$, this says
$x\in A\setminus(B\cup C)$. So
$(A\setminus B)\cap(A\setminus C)\subseteq A\setminus(B\cup C)$.

\emph{Step 7 (conclude).} Both inclusions hold, so the sets are equal.
\end{proof}
\noindent\key{Why no cases here?} There is no ``or'' among the
\emph{assumptions}. The only ``or'' is inside $x\notin B\cup C$, and
negating it turns it into an ``and''. In Practice Problem 4 the ``or'' is
an assumption, so you must split into cases.
\end{ex}

\subsection{Induction (HW1 P2, P3)}
\begin{method}
\begin{enumerate}[nosep]
  \item \textbf{Base case:} check $P(n_0)$ directly.
  \item \textbf{Write out $P(n+1)$} by replacing every $n$ with $n+1$.
        This is your target, so write it down first.
  \item \textbf{Inductive step:} assume $P(n)$. Find the piece of
        $P(n+1)$ that is exactly $P(n)$ and substitute the hypothesis.
  \item Finish with algebra.
\end{enumerate}
If the \emph{starting} index of a sum also moves with $n$, add and
subtract the boundary terms that don't match (HW1 P3 trick).
\end{method}

\begin{ex}[Bernoulli's inequality, used in Week 4]
For $x\ge -1$ and $n\in\N$: $(1+x)^n\ge 1+nx$.
\begin{proof}
Base case $n=1$: $1+x\ge 1+x$. \checkmark

Inductive step: assume $(1+x)^n\ge 1+nx$. Since $1+x\ge 0$, we may multiply
both sides by $1+x$ without flipping the inequality:
\[
(1+x)^{n+1}\ge(1+nx)(1+x)=1+(n+1)x+nx^2\ge 1+(n+1)x,
\]
because $nx^2\ge 0$.
\end{proof}

\begin{proof}[Full proof, step by step]
Fix $x\ge-1$. Let $P(n)$ be the statement $(1+x)^n\ge1+nx$. We prove $P(n)$
for all $n\in\N$ by induction.

\emph{Step 1 (base case $n=1$).} LHS $=(1+x)^1=1+x$ and RHS $=1+1\cdot x=1+x$.
Since $1+x\ge1+x$, $P(1)$ holds. \checkmark

\emph{Step 2 (write the target $P(n+1)$).}
\[
(1+x)^{n+1}\ge1+(n+1)x.
\]

\emph{Step 3 (peel off one factor).} Assume $P(n)$. Then
$(1+x)^{n+1}=(1+x)^n\cdot(1+x)$.

\emph{Step 4 (use the hypothesis $x\ge-1$).} Since $x\ge-1$, we have
$1+x\ge0$. Multiplying both sides of $P(n)$ by the \emph{nonnegative}
number $1+x$ keeps the direction of the inequality:
\[
(1+x)^n(1+x)\ge(1+nx)(1+x).
\]

\emph{Step 5 (expand).}
\[
(1+nx)(1+x)=1+x+nx+nx^2=1+(n+1)x+nx^2.
\]

\emph{Step 6 (drop the extra term).} Since $n\ge1$ and $x^2\ge0$, we have
$nx^2\ge0$, so $1+(n+1)x+nx^2\ge1+(n+1)x$.

\emph{Step 7 (chain and conclude).} Combining Steps 3--6,
\[
(1+x)^{n+1}=(1+x)^n(1+x)\ge(1+nx)(1+x)=1+(n+1)x+nx^2\ge1+(n+1)x,
\]
which is $P(n+1)$. By induction, $P(n)$ holds for all $n\in\N$.
\end{proof}
\noindent\key{Note:} the hypothesis $x\ge -1$ is used exactly once, to
multiply an inequality by $1+x\ge0$. Look for where each hypothesis is
used.

\noindent\key{What goes wrong for $x<-1$?} Take $x=-4$, $n=3$:
$(1-4)^3=-27$ but $1+3(-4)=-11$, and $-27<-11$. So the claim really can
fail when $x<-1$, and Step 4 is exactly where the proof breaks.
\end{ex}

\begin{ex}[Product identity, HW1 P2 style]
For $n\ge 2$: $\displaystyle\prod_{k=2}^{n}\Bigl(1-\frac1{k^2}\Bigr)=\frac{n+1}{2n}$.

\noindent\emph{Key step:}
$\frac{n+1}{2n}\cdot\frac{(n+1)^2-1}{(n+1)^2}
=\frac{n+1}{2n}\cdot\frac{n(n+2)}{(n+1)^2}=\frac{n+2}{2(n+1)}$, using
$(n+1)^2-1=n(n+2)$ (difference of squares).

\begin{proof}[Full proof, step by step]
Let $P(n)$ be the statement
$\displaystyle\prod_{k=2}^{n}\Bigl(1-\frac1{k^2}\Bigr)=\frac{n+1}{2n}$.
We prove $P(n)$ for all $n\ge2$ by induction.

\emph{Step 1 (base case $n=2$).} The product has a single factor,
$k=2$:
\[
\text{LHS}=1-\frac1{2^2}=1-\frac14=\frac34,\qquad
\text{RHS}=\frac{2+1}{2\cdot2}=\frac34. \checkmark
\]

\emph{Step 2 (write the target $P(n+1)$).} Replace every $n$ by $n+1$:
\[
\prod_{k=2}^{n+1}\Bigl(1-\frac1{k^2}\Bigr)=\frac{(n+1)+1}{2(n+1)}
=\frac{n+2}{2(n+1)}.
\]
This is what we must reach.

\emph{Step 3 (peel off the last factor).} Assume $P(n)$. The product up to
$n+1$ is the product up to $n$, times the $k=n+1$ factor:
\[
\prod_{k=2}^{n+1}\Bigl(1-\frac1{k^2}\Bigr)
=\underbrace{\prod_{k=2}^{n}\Bigl(1-\frac1{k^2}\Bigr)}_{\text{this is }P(n)}
\cdot\Bigl(1-\frac1{(n+1)^2}\Bigr).
\]

\emph{Step 4 (substitute the inductive hypothesis).}
\[
=\frac{n+1}{2n}\cdot\Bigl(1-\frac1{(n+1)^2}\Bigr).
\]

\emph{Step 5 (common denominator in the last factor).}
\[
1-\frac1{(n+1)^2}=\frac{(n+1)^2}{(n+1)^2}-\frac1{(n+1)^2}
=\frac{(n+1)^2-1}{(n+1)^2}.
\]

\emph{Step 6 (factor the numerator).} Use $a^2-b^2=(a-b)(a+b)$ with
$a=n+1$, $b=1$:
\[
(n+1)^2-1=\bigl((n+1)-1\bigr)\bigl((n+1)+1\bigr)=n(n+2).
\]
(Check by expanding: $(n+1)^2-1=n^2+2n+1-1=n^2+2n=n(n+2)$.)

\emph{Step 7 (multiply and cancel).}
\[
\frac{n+1}{2n}\cdot\frac{n(n+2)}{(n+1)^2}
=\frac{(n+1)\,n\,(n+2)}{2n\,(n+1)^2}
=\frac{\cancel{n}\,\cancel{(n+1)}\,(n+2)}{2\cancel{n}\,(n+1)^{\cancel{2}}}
=\frac{n+2}{2(n+1)}.
\]
Cancelling is allowed because $n\ne0$ and $n+1\ne0$ (as $n\ge2$).

\emph{Step 8 (conclude).} This is exactly the target from Step 2, so
$P(n)\Rightarrow P(n+1)$. Together with the base case, $P(n)$ holds for
all $n\ge2$.
\end{proof}

\noindent\key{Sanity check $n=3$:}
$\bigl(1-\frac14\bigr)\bigl(1-\frac19\bigr)=\frac34\cdot\frac89=\frac{24}{36}=\frac23$,
and $\frac{3+1}{2\cdot3}=\frac46=\frac23$. \checkmark

\noindent\key{Where does $\frac{n+1}{2n}$ come from? (telescoping)}
Each factor splits as
$1-\frac1{k^2}=\frac{k^2-1}{k^2}=\frac{k-1}{k}\cdot\frac{k+1}{k}$, so
\[
\prod_{k=2}^{n}\Bigl(1-\frac1{k^2}\Bigr)
=\Bigl(\frac12\cdot\frac23\cdots\frac{n-1}{n}\Bigr)
\Bigl(\frac32\cdot\frac43\cdots\frac{n+1}{n}\Bigr)
=\frac1n\cdot\frac{n+1}2=\frac{n+1}{2n}.
\]
Almost everything cancels in each bracket. This is a good way to
\emph{guess} the formula. The induction proof above is what makes it
rigorous.
\end{ex}

\subsection{Contradiction and irrationality (HW1 P4, Week 2)}
\begin{method}
Assume the opposite. Solve for the piece you know is irrational. Closure
of $\Q$ under $+,-,\times,\div$ then forces that piece to be rational,
which is a contradiction.
\end{method}

\begin{ex}[$\sqrt2\notin\Q$, from lecture: know this proof cold]
Suppose $\sqrt2=p/q$ in lowest terms. Then $p^2=2q^2$, so $p^2$ is even.
If $p=2k+1$ were odd, then $p^2=2(2k^2+2k)+1$ would be odd. So $p$ is even,
$p=2k$. Then $4k^2=2q^2$, so $q^2=2k^2$ and $q$ is even too. This
contradicts lowest terms.

\medskip
\noindent\key{Lemma used twice.} For an integer $m$: if $m^2$ is even,
then $m$ is even.

\noindent\emph{Proof (contrapositive).} Suppose $m$ is odd, so $m=2k+1$
for some $k\in\Z$. Then
\[
m^2=(2k+1)^2=4k^2+4k+1=2(2k^2+2k)+1,
\]
which is $2\cdot(\text{integer})+1$, so $m^2$ is odd. Hence ``$m$ odd
$\Rightarrow$ $m^2$ odd'', which is equivalent to ``$m^2$ even
$\Rightarrow$ $m$ even''. $\square$

\begin{proof}[Full proof, step by step]
\emph{Step 1 (assume the opposite).} Suppose, for contradiction, that
$\sqrt2\in\Q$.

\emph{Step 2 (write it in lowest terms).} Then $\sqrt2=\frac pq$ with
$p,q\in\Z$ and $q\ne0$. By cancelling common factors we may assume the
fraction is in lowest terms, i.e.\ $p$ and $q$ have \emph{no common factor
greater than $1$}. (Every fraction can be reduced, for example
$\frac{6}{4}=\frac32$.) This extra assumption is what we will contradict.

\emph{Step 3 (clear the square root).} Square both sides and multiply by
$q^2$:
\[
2=\frac{p^2}{q^2}\quad\Longrightarrow\quad p^2=2q^2.
\]

\emph{Step 4 ($p$ is even).} $p^2=2q^2$ is $2\cdot(\text{integer})$, so
$p^2$ is even. By the Lemma, $p$ is even.

\emph{Step 5 (substitute $p=2k$).} Since $p$ is even, $p=2k$ for some
$k\in\Z$. Plug into $p^2=2q^2$:
\[
(2k)^2=2q^2\quad\Longrightarrow\quad 4k^2=2q^2
\quad\Longrightarrow\quad q^2=2k^2.
\]

\emph{Step 6 ($q$ is even).} $q^2=2k^2$ is $2\cdot(\text{integer})$, so
$q^2$ is even. By the Lemma again, $q$ is even.

\emph{Step 7 (contradiction).} Now $2$ divides both $p$ and $q$, so they
share the common factor $2$. This contradicts Step 2 (lowest terms).

\emph{Step 8 (conclude).} The assumption in Step 1 is false, so
$\sqrt2\notin\Q$.
\end{proof}

\noindent\key{Why lowest terms?} Without it there is no contradiction.
Finding that $p$ and $q$ are both even would just mean the fraction could
be reduced. Lowest terms is the ``trap'' that the evenness argument springs.

\noindent\key{Common point-losers:}
\begin{itemize}[nosep]
  \item Saying ``$p^2$ even, so $p$ even'' with no justification. Prove the
        Lemma, or at least cite it.
  \item Forgetting to say $q\ne0$ and to assume lowest terms at the start.
  \item Concluding ``$p$ even'' from ``$p^2=2q^2$'' without first noting
        that $p^2$ is even.
\end{itemize}

\noindent\key{Sanity check: why doesn't this ``prove'' $\sqrt4\notin\Q$?}
Run the same steps: $p^2=4q^2$, so $p$ is even, $p=2k$. Then $4k^2=4q^2$,
so $q^2=k^2$. This says \emph{nothing} about $q$ being even, so no
contradiction arises, which is correct because $\sqrt4=2\in\Q$. The factor
$2$ appearing exactly once in $p^2=2q^2$ is what makes Step 6 work.

\noindent\key{Same template for $\sqrt3$:} replace ``even/odd'' with
``divisible by $3$ or not.'' If $3\nmid m$, then $m=3k+1$ or $m=3k+2$, and
in both cases $m^2$ leaves remainder $1$ when divided by $3$. So
$3\mid m^2\Rightarrow 3\mid m$, and the rest of the proof goes through
unchanged.
\end{ex}

\begin{ex}[rational $\pm$ irrational]
Let $r\in\Q$ and $y\notin\Q$. Then
\begin{enumerate}[nosep, label=(\alph*)]
  \item $r+y\notin\Q$ and $r-y\notin\Q$;
  \item if also $r\ne0$, then $ry\notin\Q$ and $y/r\notin\Q$.
\end{enumerate}

\noindent\key{Tool: closure of $\Q$.} If $a=\frac pq$ and $b=\frac st$ with
$p,q,s,t\in\Z$ and $q,t\ne0$, then
\[
a+b=\frac{pt+sq}{qt},\qquad a-b=\frac{pt-sq}{qt},\qquad
ab=\frac{ps}{qt},\qquad \frac ab=\frac{pt}{qs}\ \ (\text{if } b\ne0,\text{ i.e.\ } s\ne0).
\]
Each is an integer over a nonzero integer, so it is rational. This is the
only fact about $\Q$ the proof uses.

\begin{proof}[Proof of (a), the sum]
\emph{Step 1 (assume the opposite).} Suppose, for contradiction, that
$r+y\in\Q$. Give it a name: $z:=r+y$, so $z\in\Q$.

\emph{Step 2 (solve for the irrational piece).} Subtract $r$ from both
sides: $y=z-r$.

\emph{Step 3 (apply closure).} Both $z$ and $r$ are rational, so $z-r$ is
rational by closure under subtraction. Hence $y\in\Q$.

\emph{Step 4 (contradiction).} This contradicts the hypothesis
$y\notin\Q$. So the assumption in Step~1 was false, and $r+y\notin\Q$.
\end{proof}

\begin{proof}[Proof of (a), the difference]
Suppose $r-y=z\in\Q$. Solving for $y$ gives $y=r-z$, which is rational
by closure. This contradicts $y\notin\Q$, so $r-y\notin\Q$.
\end{proof}

\begin{proof}[Proof of (b), the product]
\emph{Step 1.} Suppose, for contradiction, that $ry\in\Q$. Call it
$w:=ry$, so $w\in\Q$.

\emph{Step 2.} Solve for $y$. Because $r\ne0$ we may divide by $r$:
$y=\dfrac wr$.

\emph{Step 3.} $w$ and $r$ are rational and $r\ne0$, so $w/r$ is rational
by closure under division. Hence $y\in\Q$.

\emph{Step 4.} This contradicts $y\notin\Q$, so $ry\notin\Q$.

For the quotient, note $1/r\in\Q$ and $1/r\ne0$, and $y/r=(1/r)\cdot y$.
The product case with $1/r$ in place of $r$ gives $y/r\notin\Q$.
\end{proof}

\noindent\key{Where is $r\ne0$ used?} Only in Step 2 of (b), to divide
by $r$. It is really needed: if $r=0$, then $ry=0\in\Q$, so the claim
is false.

\noindent\key{Concrete check.} $3+\sqrt2\notin\Q$: if $3+\sqrt2=z\in\Q$,
then $\sqrt2=z-3\in\Q$, which contradicts the previous example. Likewise
$5\sqrt2\notin\Q$, because otherwise $\sqrt2=\frac{5\sqrt2}{5}\in\Q$.

\noindent\key{Counterexample question:} is irrational $+$ irrational
always irrational? No: $\sqrt2+(-\sqrt2)=0$. Another one is
$\sqrt2+(1-\sqrt2)=1$, where $1-\sqrt2\notin\Q$ by part (a). Is
irrational $\times$ irrational always irrational? No: $\sqrt2\cdot\sqrt2=2$
and $\sqrt2\cdot\sqrt8=\sqrt{16}=4$.

\noindent\emph{Why the proof above fails here:} Step 3 needs $r$ to be
rational to conclude $z-r\in\Q$. If $r$ is irrational, closure says
nothing, and the argument breaks down.
\end{ex}

% =====================================================================
\section{Supremum, infimum, completeness}

\subsection{Definitions (state these exactly)}
\begin{defn}
Let $S\subseteq\R$ be nonempty.
\begin{itemize}[nosep]
  \item $M$ is an \textbf{upper bound} of $S$ if $s\le M$ for all $s\in S$.
  \item $M=\sup S$ if $M$ is an upper bound of $S$ and no $M'<M$ is an
        upper bound of $S$. (It is the \emph{least} upper bound.)
  \item $m=\inf S$ if $m$ is a lower bound of $S$ and no $m'>m$ is a lower
        bound of $S$.
  \item $\max S$ is an upper bound that \emph{belongs to} $S$. If
        $\max S$ exists, then $\sup S=\max S$.
\end{itemize}
\end{defn}

\begin{thm}[Completeness Axiom]
Every nonempty subset of $\R$ that is bounded above has a supremum in
$\R$. (It follows that every nonempty set bounded below has an infimum.)
\end{thm}

\begin{thm}[$\eps$-characterization of sup, the workhorse]
\label{thm:epssup}
$M=\sup S$ if and only if
\begin{enumerate}[nosep,label=(\roman*)]
  \item $s\le M$ for all $s\in S$, and
  \item for every $\eps>0$ there is $s\in S$ with $s>M-\eps$.
\end{enumerate}
The same holds for inf with the inequalities flipped: (ii) becomes
$\forall\eps>0\ \exists s\in S$ with $s<m+\eps$.
\end{thm}

\noindent\key{Two facts used constantly:}
\begin{itemize}[nosep]
  \item If $t$ is \emph{any} upper bound of $S$, then $\sup S\le t$.
        (This is the ``least'' part.)
  \item $\sup S-\eps$ is \emph{not} an upper bound, so some element of $S$
        exceeds it.
\end{itemize}

\subsection{Archimedean property and density}
\begin{thm}[Archimedean property]
For any $x\in\R$ there is $n\in\N$ with $n>x$. Equivalently: for any
$\eps>0$ there is $n\in\N$ with $1/n<\eps$.
\end{thm}
\noindent\emph{Proof idea (contradiction plus completeness):} if $\N$
were bounded above, let $M=\sup\N$. Then $M-1$ is not an upper bound, so
some $n>M-1$. But then $n+1>M$ with $n+1\in\N$, contradiction.

\begin{thm}[Density]
If $x<y$, there is $r\in\Q$ with $x<r<y$, and there is an irrational $z$
with $x<z<y$.
\end{thm}
\noindent\emph{Proof idea for $\Q$:} pick $n$ with $1/n<y-x$. Let
$m=\lfloor nx\rfloor+1$. Then $nx<m\le nx+1<ny$, so $x<m/n<y$.

\noindent\emph{Proof idea for irrationals (HW2 P1):} either take rationals
$x<r_1<r_2<y$ and $z=r_1+\frac{r_2-r_1}{\sqrt2}$, or apply density of $\Q$
to $x/\sqrt2<r<y/\sqrt2$ with $r\ne0$ and take $z=r\sqrt2$.

\subsection{Computing sup and inf of explicit sets}
\begin{method}
\begin{enumerate}[nosep]
  \item List the first few elements to guess the sup and inf.
  \item If the candidate is \emph{attained} (it is an element of $S$),
        show it is a bound. It is then the max or min, and you are done.
  \item If it is \emph{not attained}, show it is a bound, then use
        Theorem~\ref{thm:epssup}(ii) together with the Archimedean property.
\end{enumerate}
\end{method}

\begin{ex}
$S=\{1-\frac1n:n\in\N\}=\{0,\tfrac12,\tfrac23,\dots\}$. Then $\inf S=\min S=0$
(at $n=1$) and $\sup S=1$ (not attained).
\begin{proof}[Proof that $\sup S=1$]
(i) $1-\frac1n<1$ for all $n$, so $1$ is an upper bound. (ii) Let $\eps>0$.
By the Archimedean property, choose $n$ with $\frac1n<\eps$. Then
$1-\frac1n>1-\eps$, and $1-\frac1n\in S$. By Theorem~\ref{thm:epssup},
$\sup S=1$.
\end{proof}

\begin{proof}[Full proof, step by step]
\emph{Step 1 ($\inf S=\min S=0$).} For every $n\in\N$, $n\ge1$, so
$\frac1n\le1$ and $1-\frac1n\ge0$. Thus $0$ is a lower bound of $S$. Also
$0=1-\frac11\in S$. A lower bound that belongs to the set is the minimum,
so $\inf S=\min S=0$.

\emph{Step 2 ($1$ is an upper bound).} For every $n\in\N$, $\frac1n>0$, so
$1-\frac1n<1$. Hence $s\le1$ for all $s\in S$. This is part (i) of
Theorem~\ref{thm:epssup}.

\emph{Step 3 (start part (ii)).} Let $\eps>0$ be arbitrary. We must find
an element of $S$ that is larger than $1-\eps$.

\emph{Step 4 (scratch: what $n$ works?).}
$1-\frac1n>1-\eps\iff\frac1n<\eps$. So we need $n$ with $\frac1n<\eps$.

\emph{Step 5 (Archimedean property).} Such an $n\in\N$ exists by the
Archimedean property. Fix one.

\emph{Step 6 (verify).} Then $s:=1-\frac1n\in S$ and
$s=1-\frac1n>1-\eps$.

\emph{Step 7 (conclude).} Steps 2--6 verify (i) and (ii) of
Theorem~\ref{thm:epssup}, so $\sup S=1$.

\emph{Step 8 (not attained).} If $1\in S$ then $1=1-\frac1n$ for some $n$,
so $\frac1n=0$, which is impossible. So $1\notin S$ and $S$ has no maximum.
\end{proof}
\end{ex}

\begin{ex}[HW2 P2]
$S_1=\{\frac1n\}$: $\sup=\max=1$, and $\inf=0$ is not attained.
$S_2=\{\frac{(-1)^n}{n}\}$: split into odd and even $n$. Then
$\sup S_2=\max=\frac12$ (at $n=2$) and $\inf S_2=\min=-1$ (at $n=1$).

\begin{proof}[Full proof for $S_1$, step by step]
\emph{Step 1 ($\sup S_1=\max S_1=1$).} For every $n\in\N$, $n\ge1$, so
$\frac1n\le1$. So $1$ is an upper bound, and $1=\frac11\in S_1$. An upper
bound that belongs to the set is the maximum, so $\sup S_1=\max S_1=1$.

\emph{Step 2 ($0$ is a lower bound).} For every $n$, $\frac1n>0$. So
$0\le s$ for all $s\in S_1$.

\emph{Step 3 (the $\eps$ part of the inf).} Let $\eps>0$. By the
Archimedean property there is $n\in\N$ with $\frac1n<\eps=0+\eps$, and
$\frac1n\in S_1$. By Theorem~\ref{thm:epssup} (inf version), $\inf S_1=0$.

\emph{Step 4 (not attained).} $\frac1n=0$ has no solution, so
$0\notin S_1$ and $S_1$ has no minimum.
\end{proof}

\begin{proof}[Full proof for $S_2$, step by step]
\emph{Step 1 (list and split).} The elements are
$-1,\frac12,-\frac13,\frac14,-\frac15,\dots$.
\begin{itemize}[nosep]
  \item If $n$ is even, then $\frac{(-1)^n}{n}=\frac1n$, and $n\ge2$, so
        $0<\frac1n\le\frac12$.
  \item If $n$ is odd, then $\frac{(-1)^n}{n}=-\frac1n$, and $n\ge1$, so
        $-1\le-\frac1n<0$.
\end{itemize}

\emph{Step 2 (upper bound $\frac12$).} Even terms are $\le\frac12$ by
Step 1. Odd terms are $<0<\frac12$. So every element is $\le\frac12$.

\emph{Step 3 (attained).} $n=2$ gives $\frac{(-1)^2}{2}=\frac12\in S_2$.
So $\sup S_2=\max S_2=\frac12$.

\emph{Step 4 (lower bound $-1$).} Odd terms are $\ge-1$ by Step 1. Even
terms are $>0>-1$. So every element is $\ge-1$.

\emph{Step 5 (attained).} $n=1$ gives $-1\in S_2$. So
$\inf S_2=\min S_2=-1$.
\end{proof}
\noindent\key{Lesson:} when the candidate is attained, you never need
$\eps$. ``Bound $+$ element of the set'' finishes it.
\end{ex}

\begin{ex}
$S=(a,b)$: $\sup S=b$ and $\inf S=a$, and neither is attained. For (ii):
given $\eps>0$ (with $\eps<b-a$ if needed), the point
$x=b-\eps/2\in S$ satisfies $x>b-\eps$.

\begin{proof}[Full proof that $\sup(a,b)=b$ (for $a<b$), step by step]
\emph{Step 1 (upper bound).} If $x\in(a,b)$ then $a<x<b$, so $x\le b$.
Thus $b$ is an upper bound.

\emph{Step 2 (start part (ii)).} Let $\eps>0$ be arbitrary. We need
$x\in(a,b)$ with $x>b-\eps$.

\emph{Step 3 (handle large $\eps$).} The point $b-\frac\eps2$ might fall
at or below $a$ when $\eps$ is large. To avoid this, let
$\delta=\min\{\eps,\ b-a\}>0$, and set $x=b-\frac\delta2$.

\emph{Step 4 ($x\in(a,b)$).} Since $\delta>0$, $x<b$. Since
$\delta\le b-a$, we have $\frac\delta2<b-a$, so $x=b-\frac\delta2>b-(b-a)=a$.
So $a<x<b$.

\emph{Step 5 ($x>b-\eps$).} Since $\delta\le\eps$, we have
$\frac\delta2<\eps$, so $x=b-\frac\delta2>b-\eps$.

\emph{Step 6 (conclude).} By Theorem~\ref{thm:epssup}, $\sup(a,b)=b$.

\emph{Step 7 (not attained).} $b\notin(a,b)$ since $b<b$ is false. So
$(a,b)$ has no maximum.
\end{proof}
The proof that $\inf(a,b)=a$ is the mirror image: use
$x=a+\frac\delta2$ with the same $\delta$.
\end{ex}

\subsection{Standard sup/inf proofs: templates}

\begin{ex}[$\sup S\le\inf T$, HW2 P3]
If $s\le t$ for all $s\in S$, $t\in T$: fix $t$. Then $t$ is an upper bound
of $S$, so $\sup S\le t$. Since $t$ was arbitrary, $\sup S$ is a lower
bound of $T$, so $\sup S\le\inf T$.

\begin{proof}[Full proof, step by step]
Let $S,T\subseteq\R$ be nonempty with $s\le t$ for all $s\in S$, $t\in T$.

\emph{Step 1 ($\sup S$ exists).} $T\ne\emptyset$, so pick some $t_0\in T$.
By hypothesis $s\le t_0$ for all $s\in S$, so $S$ is bounded above. Since
$S\ne\emptyset$, the Completeness Axiom gives that $\sup S$ exists.
Similarly, any $s_0\in S$ is a lower bound of $T$, so $\inf T$ exists.

\emph{Step 2 (fix an arbitrary $t$).} Let $t\in T$ be arbitrary.

\emph{Step 3 ($t$ is an upper bound of $S$).} By hypothesis $s\le t$ for
every $s\in S$. That is the definition of ``$t$ is an upper bound of $S$''.

\emph{Step 4 (``least'' upper bound).} $\sup S$ is less than or equal to
\emph{every} upper bound of $S$. So $\sup S\le t$.

\emph{Step 5 (turn it around).} Since $t\in T$ was arbitrary, Step 4 says
$\sup S\le t$ for all $t\in T$. So $\sup S$ is a lower bound of $T$.

\emph{Step 6 (``greatest'' lower bound).} $\inf T$ is greater than or
equal to every lower bound of $T$. So $\sup S\le\inf T$.
\end{proof}
\noindent\key{Pattern:} ``fix one, bound, then let it vary'' is used
twice. You apply the ``least'' property to $S$ and then the ``greatest''
property to $T$.
\end{ex}

\begin{ex}[$\sup(A+B)=\sup A+\sup B$, HW2 P4]
($\le$) Since $a+b\le\sup A+\sup B$, the right side is an upper bound.
($\ge$) Given $\eps>0$, pick $a>\sup A-\frac\eps2$ and
$b>\sup B-\frac\eps2$. Then $a+b>\sup A+\sup B-\eps$.
\key{Splitting $\eps$ into halves} is the standard trick.

\begin{proof}[Full proof, step by step]
Let $A,B$ be nonempty and bounded above, and
$A+B=\{a+b:a\in A,\ b\in B\}$. Write $\alpha=\sup A$ and $\beta=\sup B$
(these exist by completeness). We show $\sup(A+B)=\alpha+\beta$ using
Theorem~\ref{thm:epssup}.

\emph{Step 1 ($A+B$ is nonempty).} Pick $a_0\in A$ and $b_0\in B$. Then
$a_0+b_0\in A+B$.

\emph{Step 2 (upper bound).} Let $x\in A+B$ be arbitrary. Then $x=a+b$ for
some $a\in A$, $b\in B$. Since $a\le\alpha$ and $b\le\beta$, adding gives
$x=a+b\le\alpha+\beta$. So $\alpha+\beta$ is an upper bound of $A+B$. In
particular $\sup(A+B)$ exists.

\emph{Step 3 (start part (ii)).} Let $\eps>0$. We need some $x\in A+B$
with $x>\alpha+\beta-\eps$.

\emph{Step 4 (split $\eps$).} $\frac\eps2>0$, so $\alpha-\frac\eps2$ is not
an upper bound of $A$. Hence there is $a\in A$ with $a>\alpha-\frac\eps2$.
In the same way there is $b\in B$ with $b>\beta-\frac\eps2$.

\emph{Step 5 (add).} Then $a+b\in A+B$ and
\[
a+b>\Bigl(\alpha-\frac\eps2\Bigr)+\Bigl(\beta-\frac\eps2\Bigr)=\alpha+\beta-\eps.
\]

\emph{Step 6 (conclude).} Steps 2 and 5 are (i) and (ii) of
Theorem~\ref{thm:epssup}, so $\sup(A+B)=\alpha+\beta$.
\end{proof}
\noindent\key{Why $\frac\eps2$ and not $\eps$?} With $a>\alpha-\eps$ and
$b>\beta-\eps$ you only get $a+b>\alpha+\beta-2\eps$, which is too weak.
Each of the two pieces gets half of the budget.
\end{ex}

\begin{ex}[$\sup(cA)=c\sup A$ for $c>0$]
Let $\alpha=\sup A$. For $a\in A$, $ca\le c\alpha$, so $c\alpha$ is an
upper bound of $cA$. Given $\eps>0$, pick $a\in A$ with
$a>\alpha-\eps/c$. Then $ca>c\alpha-\eps$. \emph{(What changes if $c<0$?
Then $\sup(cA)=c\inf A$.)}

\begin{proof}[Full proof, step by step]
Let $A$ be nonempty and bounded above, $c>0$, and
$cA=\{ca:a\in A\}$. Let $\alpha=\sup A$.

\emph{Step 1 (upper bound).} Let $x\in cA$, so $x=ca$ for some $a\in A$.
Since $a\le\alpha$ and $c>0$, multiplying by $c$ keeps the inequality:
$ca\le c\alpha$. So $c\alpha$ is an upper bound of $cA$.

\emph{Step 2 (start part (ii)).} Let $\eps>0$. We want $x\in cA$ with
$x>c\alpha-\eps$.

\emph{Step 3 (scratch: rescale $\eps$).} If $x=ca$, then
$ca>c\alpha-\eps\iff a>\alpha-\frac\eps c$ (divide by $c>0$). So we need
an element of $A$ within $\frac\eps c$ of $\alpha$.

\emph{Step 4 (use $\sup A$).} $\frac\eps c>0$, so
$\alpha-\frac\eps c$ is not an upper bound of $A$. So there is $a\in A$
with $a>\alpha-\frac\eps c$.

\emph{Step 5 (multiply back).} Multiply by $c>0$:
$ca>c\alpha-\eps$, and $ca\in cA$.

\emph{Step 6 (conclude).} By Theorem~\ref{thm:epssup},
$\sup(cA)=c\alpha=c\sup A$.
\end{proof}
\noindent\key{Where is $c>0$ used?} In Steps 1, 3 and 5, every time we
multiply or divide an inequality by $c$. If $c<0$ each of those
inequalities flips, upper bounds turn into lower bounds, and you get
$\sup(cA)=c\inf A$ instead.
\end{ex}

\begin{ex}[Quiz 1: $\sup(AB)=\sup A\cdot\sup B$ for positive sets]\label{ex:quiz1}
Let $A,B\subseteq(0,\infty)$ be nonempty and bounded, with $\alpha=\sup A$,
$\beta=\sup B$, and $C=\{xy:x\in A,\ y\in B\}$. Prove $\sup C=\alpha\beta$.

\smallskip\noindent\textbf{Idea.} ($\le$) is the usual ``upper bound''
step. For ($\ge$), the instructor's solution avoids $\eps$: it
\key{divides by one variable to turn the problem into a sup of $A$}, then
does the same for $B$. This is ``fix one, bound, let it vary'' from the
$\sup S\le\inf T$ template, done twice.

\begin{proof}[Full proof, step by step]
\emph{Step 1 ($\alpha\beta$ is an upper bound of $C$).} Let $z\in C$, so
$z=xy$ with $x\in A$, $y\in B$. Then $0<x\le\alpha$ and $0<y\le\beta$.
Multiplying inequalities between \emph{positive} numbers,
$xy\le\alpha y\le\alpha\beta$. So $\alpha\beta$ is an upper bound of $C$,
$\sup C$ exists by completeness ($C\ne\emptyset$), and
$\sup C\le\alpha\beta$.

\emph{Step 2 (fix $y$, bound $A$).} Fix $y\in B$. For every $x\in A$,
$xy\in C$, so $xy\le\sup C$. Since $y>0$, $x\le\frac{\sup C}{y}$. So
$\frac{\sup C}{y}$ is an upper bound of $A$, and since $\alpha$ is the
\emph{least} upper bound, $\alpha\le\frac{\sup C}{y}$.

\emph{Step 3 (let $y$ vary, bound $B$).} Step 2 gives
$y\le\frac{\sup C}{\alpha}$ for every $y\in B$ (here $\alpha>0$ because
$A$ contains a positive number). So $\frac{\sup C}{\alpha}$ is an upper
bound of $B$, and therefore $\beta\le\frac{\sup C}{\alpha}$, i.e.\
$\alpha\beta\le\sup C$.

\emph{Step 4.} Steps 1 and 3 give $\sup C=\alpha\beta$.
\end{proof}
\noindent\key{Where is positivity used?} In every multiplication or
division (Steps 1, 2, 3). Without it the result fails: for
$A=B=\{-1,-2\}$ we get $\sup A=\sup B=-1$, but $C=\{1,2,4\}$ and
$\sup C=4\ne1$.

\noindent\key{$\eps$ version (harder).} You can also pick $x>\alpha-\delta$,
$y>\beta-\delta$ and expand
$(\alpha-\delta)(\beta-\delta)>\alpha\beta-\delta(\alpha+\beta)$, but you
then have to choose $\delta$ small enough to keep both factors positive.
The division trick above is cleaner. \textbf{Expect a variation on the
exam}, e.g.\ $\sup\{1/x:x\in A\}=1/\inf A$ (Practice exam,
Problem~\ref{mock:recip}).
\end{ex}

\begin{ex}[$\inf S=-\sup(-S)$]
$m\le s$ for all $s\in S$ $\iff$ $-s\le -m$ for all $s\in S$. So lower
bounds of $S$ are exactly the negatives of upper bounds of $-S$, and the
greatest lower bound corresponds to the least upper bound.

\begin{proof}[Full proof, step by step]
Let $S$ be nonempty and bounded below, and let $-S=\{-s:s\in S\}$.

\emph{Step 1 (key equivalence).} For any real $m$:
\begin{align*}
m\text{ is a lower bound of }S
&\iff m\le s\ \text{ for all } s\in S\\
&\iff -s\le -m\ \text{ for all } s\in S\\
&\iff -m\text{ is an upper bound of }-S.
\end{align*}
(Multiplying by $-1$ flips inequalities.)

\emph{Step 2 ($\sup(-S)$ exists).} $S$ has some lower bound $m$, so by
Step 1 $-m$ is an upper bound of $-S$. Also $-S\ne\emptyset$. By
completeness, $\beta:=\sup(-S)$ exists.

\emph{Step 3 ($-\beta$ is a lower bound of $S$).} $\beta$ is an upper
bound of $-S$, i.e.\ $-(-\beta)$ is an upper bound of $-S$. By Step 1
(with $m=-\beta$), $-\beta$ is a lower bound of $S$.

\emph{Step 4 ($-\beta$ is the greatest).} Let $m$ be any lower bound of
$S$. By Step 1, $-m$ is an upper bound of $-S$. Since $\beta$ is the
\emph{least} upper bound, $\beta\le-m$, so $m\le-\beta$.

\emph{Step 5 (conclude).} $-\beta$ is a lower bound of $S$ and every lower
bound is $\le-\beta$. By definition, $\inf S=-\beta=-\sup(-S)$.
\end{proof}
\noindent\key{Use:} this lets you get every inf fact from the matching sup
fact, and also shows that completeness for sups gives completeness for
infs.
\end{ex}

\subsection{Monotonicity of sup}
If $A\subseteq B$, then $\sup A\le\sup B$ and $\inf A\ge\inf B$.
(Any upper bound of $B$ is an upper bound of $A$.)

% =====================================================================
\section{Absolute value}

\begin{thm}[Properties]\
\begin{enumerate}[nosep]
  \item $|xy|=|x||y|$ and $|-x|=|x|$.
  \item For $a\ge0$: $|x|\le a\iff -a\le x\le a$. \key{(Most useful fact.)}
  \item Triangle inequality: $|x+y|\le|x|+|y|$.
  \item Reverse triangle inequality: $\bigl||x|-|y|\bigr|\le|x-y|$.
  \item $|x-a|<\delta\iff a-\delta<x<a+\delta$.
\end{enumerate}
\end{thm}

\begin{method}[Proving the triangle inequality]
Add $-|x|\le x\le|x|$ and $-|y|\le y\le|y|$ to get
$-(|x|+|y|)\le x+y\le|x|+|y|$, then apply property 2.
\end{method}

\begin{method}[Reverse triangle inequality, HW3 P1]
Write $x=(x-y)+y$, so $|x|\le|x-y|+|y|$, i.e.\ $|x|-|y|\le|x-y|$. Swap
$x$ and $y$ to get $-(|x|-|y|)\le|x-y|$. Then property 2 gives the result.
\end{method}

\begin{method}[The add-and-subtract trick]
To compare $a$ and $b$ through $c$, write $a-b=(a-c)+(c-b)$ and apply the
triangle inequality. This is how uniqueness of limits and the product rule
are proved.
\end{method}

\begin{ex}[Solving equations with absolute values]
Solve $|x-3|+|x-1|=3$. Break the line at the points $1$ and $3$:
\begin{itemize}[nosep]
  \item $x<1$: $(3-x)+(1-x)=3\Rightarrow x=\frac12$. \checkmark
  \item $1\le x<3$: $(3-x)+(x-1)=2\ne3$. No solution.
  \item $x\ge3$: $(x-3)+(x-1)=3\Rightarrow x=\frac72$. \checkmark
\end{itemize}

\begin{proof}[Full solution, step by step]
\emph{Step 1 (find the break points).} $|x-3|$ changes formula at $x=3$,
and $|x-1|$ changes formula at $x=1$. These points split $\R$ into
$(-\infty,1)$, $[1,3)$ and $[3,\infty)$.

\emph{Step 2 (sign table).} Use $|u|=u$ if $u\ge0$ and $|u|=-u$ if $u<0$:
\begin{center}
\begin{tabular}{c|c|c}
 & $x-3$ & $x-1$\\ \hline
$x<1$ & $<0$, so $|x-3|=3-x$ & $<0$, so $|x-1|=1-x$\\
$1\le x<3$ & $<0$, so $|x-3|=3-x$ & $\ge0$, so $|x-1|=x-1$\\
$x\ge3$ & $\ge0$, so $|x-3|=x-3$ & $>0$, so $|x-1|=x-1$
\end{tabular}
\end{center}

\emph{Step 3 (case $x<1$).} The equation becomes $(3-x)+(1-x)=3$, i.e.\
$4-2x=3$, so $x=\frac12$. Check that it lies in the case: $\frac12<1$.
\checkmark

\emph{Step 4 (case $1\le x<3$).} The equation becomes $(3-x)+(x-1)=3$,
i.e.\ $2=3$. This is false, so there is no solution in this interval.

\emph{Step 5 (case $x\ge3$).} The equation becomes $(x-3)+(x-1)=3$, i.e.\
$2x-4=3$, so $x=\frac72$. Check: $\frac72\ge3$. \checkmark

\emph{Step 6 (verify in the original equation).}
$|\tfrac12-3|+|\tfrac12-1|=\tfrac52+\tfrac12=3$ and
$|\tfrac72-3|+|\tfrac72-1|=\tfrac12+\tfrac52=3$.

\emph{Step 7 (answer).} The solution set is $\{\frac12,\frac72\}$.
\end{proof}
\noindent\key{Always check that a case's answer lies in that case's
interval.} If it doesn't, throw it away.

\noindent\key{Geometric check:} $|x-3|+|x-1|$ is the total distance from
$x$ to $1$ and to $3$. Between $1$ and $3$ this total is always $2$. You
must move $\frac12$ outside the interval on either side to get $3$.
\end{ex}

% =====================================================================
\section{Limits of sequences: the $\eps$--$N$ definition}

\begin{defn}
$x_n\to L$ means: for every $\eps>0$ there exists $N\in\N$ such that
$|x_n-L|<\eps$ for all $n\ge N$.
\end{defn}

\begin{method}[The $\eps$--$N$ recipe]
\begin{enumerate}[nosep]
  \item \textbf{Guess $L$} (leading terms, or intuition from calculus).
  \item \textbf{Scratch work:} simplify $|x_n-L|$ to a single fraction.
  \item \textbf{Overestimate} it by something simpler, such as $\frac{C}{n}$.
        You make a fraction larger by making the denominator smaller or the
        numerator larger. You only need an upper bound.
  \item Solve $\frac{C}{n}<\eps$, which gives $n>\frac{C}{\eps}$, so
        $N>\frac{C}{\eps}$ (Archimedean).
  \item \textbf{Write it forwards:} ``Let $\eps>0$. Choose $N>\ldots$.
        Then for $n\ge N$: $|x_n-L|=\cdots\le\cdots<\eps$.''
\end{enumerate}
\end{method}

\begin{ex}[HW3 P2]
$\frac{n+2}{3n+5}\to\frac13$.
\begin{scratch}
$\left|\frac{n+2}{3n+5}-\frac13\right|=\frac{1}{3(3n+5)}=\frac1{9n+15}<\frac1{9n}$.
We want $\frac1{9n}<\eps$, i.e.\ $n>\frac1{9\eps}$.
\end{scratch}
\begin{proof}
Let $\eps>0$ and choose $N\in\N$ with $N>\frac{1}{9\eps}$. For $n\ge N$,
$\left|\frac{n+2}{3n+5}-\frac13\right|=\frac1{9n+15}<\frac1{9n}\le\frac1{9N}<\eps$.
\end{proof}

\begin{proof}[Full proof, step by step]
\emph{Step 1 (simplify $|x_n-L|$, scratch).} Use the common denominator
$3(3n+5)$:
\[
\frac{n+2}{3n+5}-\frac13=\frac{3(n+2)-(3n+5)}{3(3n+5)}=\frac{3n+6-3n-5}{9n+15}
=\frac1{9n+15}.
\]
This is positive, so $\left|\frac{n+2}{3n+5}-\frac13\right|=\frac1{9n+15}$.

\emph{Step 2 (overestimate, scratch).} Since $9n+15>9n>0$,
$\frac1{9n+15}<\frac1{9n}$. (A smaller denominator gives a larger fraction.)

\emph{Step 3 (solve for $n$, scratch).}
$\frac1{9n}<\eps\iff9n\eps>1\iff n>\frac1{9\eps}$.

\emph{Step 4 (the proof starts here: let $\eps>0$).} Let $\eps>0$ be
arbitrary.

\emph{Step 5 (choose $N$).} By the Archimedean property, choose
$N\in\N$ with $N>\frac1{9\eps}$. Note that $N$ depends only on $\eps$.

\emph{Step 6 (verify for every $n\ge N$).} Let $n\ge N$. Then
$9n\ge9N>\frac1\eps$, so $\frac1{9n}\le\frac1{9N}<\eps$. Therefore
\[
\left|\frac{n+2}{3n+5}-\frac13\right|=\frac1{9n+15}<\frac1{9n}\le\frac1{9N}<\eps.
\]

\emph{Step 7 (conclude).} Since $\eps>0$ was arbitrary,
$\frac{n+2}{3n+5}\to\frac13$ by definition.
\end{proof}
\end{ex}

\begin{ex}[Lecture]
$\frac{2n-1}{n+1}\to2$: $\left|\frac{2n-1}{n+1}-2\right|=\frac{3}{n+1}<\frac3n$.
Choose $N>\frac3\eps$.

\begin{proof}[Full proof, step by step]
\emph{Step 1 (simplify, scratch).}
\[
\frac{2n-1}{n+1}-2=\frac{2n-1-2(n+1)}{n+1}=\frac{-3}{n+1},
\qquad\text{so}\qquad
\left|\frac{2n-1}{n+1}-2\right|=\frac3{n+1}.
\]

\emph{Step 2 (overestimate, scratch).} $n+1>n$, so $\frac3{n+1}<\frac3n$.
Then $\frac3n<\eps\iff n>\frac3\eps$.

\emph{Step 3 (let $\eps>0$ and choose $N$).} Let $\eps>0$. By the
Archimedean property choose $N\in\N$ with $N>\frac3\eps$.

\emph{Step 4 (verify).} For $n\ge N$:
\[
\left|\frac{2n-1}{n+1}-2\right|=\frac3{n+1}<\frac3n\le\frac3N<\eps.
\]

\emph{Step 5 (conclude).} So $\frac{2n-1}{n+1}\to2$.
\end{proof}
\noindent\key{Watch the sign:} the difference is $\frac{-3}{n+1}$, which
is negative. The absolute value removes the minus sign, so don't skip it.
\end{ex}

\begin{ex}[A numerator that needs bounding]
$\frac{3n^2+n}{n^2+1}\to3$:
$\left|\frac{3n^2+n}{n^2+1}-3\right|=\frac{|n-3|}{n^2+1}\le\frac{n+3}{n^2}\le\frac{4n}{n^2}=\frac4n$
for $n\ge1$ (since $3\le 3n$). Choose $N>\frac4\eps$.
\key{Bound the numerator above and the denominator below.}

\begin{proof}[Full proof, step by step]
\emph{Step 1 (simplify, scratch).}
\[
\frac{3n^2+n}{n^2+1}-3=\frac{3n^2+n-3n^2-3}{n^2+1}=\frac{n-3}{n^2+1},
\qquad\text{so}\qquad
\left|\frac{3n^2+n}{n^2+1}-3\right|=\frac{|n-3|}{n^2+1}.
\]

\emph{Step 2 (numerator: bound above).} By the triangle inequality,
$|n-3|\le|n|+|-3|=n+3$. Since $n\ge1$, $3\le3n$, so $n+3\le4n$.

\emph{Step 3 (denominator: bound below).} $n^2+1>n^2$, so
$\frac1{n^2+1}<\frac1{n^2}$.

\emph{Step 4 (combine).}
\[
\frac{|n-3|}{n^2+1}<\frac{4n}{n^2}=\frac4n.
\]
(For $n=3$ the left side is $0$, which is still $<\frac4n$.)

\emph{Step 5 (let $\eps>0$ and choose $N$).} Let $\eps>0$ and choose
$N\in\N$ with $N>\frac4\eps$.

\emph{Step 6 (verify).} For $n\ge N$:
$\left|\frac{3n^2+n}{n^2+1}-3\right|<\frac4n\le\frac4N<\eps$.

\emph{Step 7 (conclude).} So $\frac{3n^2+n}{n^2+1}\to3$.
\end{proof}
\noindent\key{Why not solve $\frac{|n-3|}{n^2+1}<\eps$ exactly?} You could,
but it needs the quadratic formula and cases. Overestimating by $\frac4n$
is valid and much faster. You only need \emph{some} $N$ that works, not
the smallest one.
\end{ex}

\begin{ex}[Conjugate trick, lecture]
$\frac{\sqrt{n^2+n}}{n}\to1$:
$\frac{\sqrt{n^2+n}-n}{n}=\frac{n}{n(\sqrt{n^2+n}+n)}=\frac1{\sqrt{n^2+n}+n}<\frac1n$.
Choose $N>\frac1\eps$. \key{Use $\sqrt a-\sqrt b=\frac{a-b}{\sqrt a+\sqrt b}$.}

\begin{proof}[Full proof, step by step]
\emph{Step 1 (common denominator, scratch).}
$\frac{\sqrt{n^2+n}}{n}-1=\frac{\sqrt{n^2+n}-n}{n}$.

\emph{Step 2 (multiply by the conjugate).}
\[
\sqrt{n^2+n}-n=\frac{(\sqrt{n^2+n}-n)(\sqrt{n^2+n}+n)}{\sqrt{n^2+n}+n}
=\frac{(n^2+n)-n^2}{\sqrt{n^2+n}+n}=\frac{n}{\sqrt{n^2+n}+n}.
\]

\emph{Step 3 (substitute and cancel).}
\[
\frac{\sqrt{n^2+n}-n}{n}=\frac{n}{n(\sqrt{n^2+n}+n)}=\frac1{\sqrt{n^2+n}+n}.
\]
This is positive, so it equals its absolute value.

\emph{Step 4 (overestimate).} $\sqrt{n^2+n}>0$, so
$\sqrt{n^2+n}+n>n$ and $\frac1{\sqrt{n^2+n}+n}<\frac1n$.

\emph{Step 5 (let $\eps>0$ and choose $N$).} Let $\eps>0$ and choose
$N\in\N$ with $N>\frac1\eps$.

\emph{Step 6 (verify).} For $n\ge N$:
\[
\left|\frac{\sqrt{n^2+n}}{n}-1\right|=\frac1{\sqrt{n^2+n}+n}<\frac1n\le\frac1N<\eps.
\]

\emph{Step 7 (conclude).} So $\frac{\sqrt{n^2+n}}{n}\to1$.
\end{proof}
\noindent\key{Why the conjugate?} $\sqrt{n^2+n}-n$ is a difference of two
large, nearly equal numbers, so you can't see its size. After multiplying
by the conjugate, the $n^2$ terms cancel.
\end{ex}

\begin{ex}[Bounded over growing]
$\frac{\sin n}{n}\to0$: $\left|\frac{\sin n}{n}\right|\le\frac1n<\eps$ when
$n>\frac1\eps$. (You cannot use the quotient rule, because $\lim\sin n$ does
not exist.)

\begin{proof}[Full proof, step by step]
\emph{Step 1 (the only fact about $\sin$).} $|\sin t|\le1$ for all real
$t$. In particular $|\sin n|\le1$.

\emph{Step 2 (bound the term).} For $n\in\N$:
$\left|\frac{\sin n}{n}-0\right|=\frac{|\sin n|}{n}\le\frac1n$.

\emph{Step 3 (let $\eps>0$ and choose $N$).} Let $\eps>0$ and choose
$N\in\N$ with $N>\frac1\eps$.

\emph{Step 4 (verify).} For $n\ge N$:
$\left|\frac{\sin n}n-0\right|\le\frac1n\le\frac1N<\eps$.

\emph{Step 5 (conclude).} So $\frac{\sin n}{n}\to0$.
\end{proof}
\end{ex}

\begin{ex}[Proving divergence]
\label{ex:diverge}
$x_n=(-1)^n$ has no limit. Suppose $x_n\to L$. Take $\eps=\frac12$ and get
$N$. Choose an even $n\ge N$ and an odd $m\ge N$. Then
$|1-L|<\frac12$ and $|-1-L|<\frac12$, so
$2=|1-(-1)|\le|1-L|+|L-(-1)|<1$, which is a contradiction.

\begin{proof}[Full proof, step by step]
\emph{Step 1 (assume the opposite).} Suppose, for contradiction, that
$(-1)^n\to L$ for some $L\in\R$.

\emph{Step 2 (choose one specific $\eps$).} The definition holds for
\emph{every} $\eps>0$, so it holds for $\eps=\frac12$. So there is
$N\in\N$ with $|(-1)^n-L|<\frac12$ for all $n\ge N$.

\emph{Step 3 (pick one even and one odd index past $N$).} Take $n=2N$
(even) and $m=2N+1$ (odd). Both are $\ge N$.

\emph{Step 4 (read off the two inequalities).} $(-1)^{2N}=1$ and
$(-1)^{2N+1}=-1$, so
\[
|1-L|<\tfrac12\qquad\text{and}\qquad|-1-L|<\tfrac12.
\]

\emph{Step 5 (add and subtract $L$, then triangle inequality).}
\[
2=|1-(-1)|=|(1-L)+(L-(-1))|\le|1-L|+|L+1|<\tfrac12+\tfrac12=1.
\]
(Here $|L+1|=|-1-L|$.)

\emph{Step 6 (contradiction).} Step 5 says $2<1$, which is false. So no
such $L$ exists, and $(-1)^n$ diverges.
\end{proof}
\noindent\key{Why $\eps=\frac12$?} The terms $1$ and $-1$ are $2$ apart.
Any $\eps\le1$ works, because two points within $\eps$ of $L$ are less
than $2\eps\le2$ apart. $\eps=\frac12$ leaves some room.
\end{ex}

\subsection{Limits equal to $\pm\infty$}
\begin{defn}
$x_n\to\infty$ means: for every $M$ there is $N$ such that $x_n>M$ for all
$n\ge N$. Also, $x_n\to-\infty$ means $-x_n\to\infty$.
\end{defn}

\begin{method}
Find a \emph{simpler lower bound} $x_n\ge g(n)$ with $g(n)\to\infty$
obviously, such as $g(n)=cn$. Then choose $N$ so that $g(n)>M$.
\end{method}

\begin{ex}[Lecture]
$\frac{n^4-4n}{n^2}\to\infty$: for $n\ge2$, $n^4-4n\ge n^3$ (since
$n^3(n-1)\ge n^3\ge 4n$ when $n\ge 2$), so $\frac{n^4-4n}{n^2}\ge n$. Given
$M$, choose $N>\max(M,2)$.

\begin{proof}[Full proof, step by step]
\emph{Step 1 (simplify).} $\frac{n^4-4n}{n^2}=n^2-\frac4n$. We'll work with
the original form, since the lower bound is easy to see there.

\emph{Step 2 (lower bound on the numerator, for $n\ge2$).} We claim
$n^4-4n\ge n^3$, i.e.\ $n^4-n^3\ge4n$. Factor:
$n^4-n^3=n^3(n-1)$. For $n\ge2$ we have $n-1\ge1$, so
$n^3(n-1)\ge n^3=n\cdot n^2\ge n\cdot4=4n$ (since $n^2\ge4$).

\emph{Step 3 (simpler lower bound for $x_n$).} Dividing Step 2 by
$n^2>0$: for $n\ge2$, $\frac{n^4-4n}{n^2}\ge\frac{n^3}{n^2}=n$.

\emph{Step 4 (let $M$ be arbitrary and choose $N$).} Let $M\in\R$. By the
Archimedean property choose $N\in\N$ with $N>\max(M,2)$.

\emph{Step 5 (verify).} Let $n\ge N$. Then $n\ge2$, so Step 3 applies, and
$\frac{n^4-4n}{n^2}\ge n\ge N>M$.

\emph{Step 6 (conclude).} For every $M$ there is such an $N$, so
$\frac{n^4-4n}{n^2}\to\infty$.
\end{proof}
\noindent\key{Why the $\max$?} $N>M$ gives the size condition, and $N\ge2$
makes Step 3 valid. Taking the $\max$ gets both at once.
\end{ex}

\begin{ex}[HW3 P3: $\infty+$bounded$=\infty$]
Say $|y_n|\le B$. Then $x_n+y_n\ge x_n-B$. Given $K$, apply
$x_n\to\infty$ with the \key{threshold $K+B$} to get $N$. For $n\ge N$,
$x_n+y_n>K+B-B=K$.

\begin{proof}[Full proof, step by step]
Assume $x_n\to\infty$ and $(y_n)$ is bounded. We show $x_n+y_n\to\infty$.

\emph{Step 1 (name the bound).} Since $(y_n)$ is bounded, there is $B>0$
with $|y_n|\le B$ for all $n$. In particular $-B\le y_n$ for all $n$.

\emph{Step 2 (lower bound for the sum).} Adding $x_n$ gives
$x_n+y_n\ge x_n-B$ for all $n$.

\emph{Step 3 (let $K$ be arbitrary).} Let $K\in\R$. We must find $N$ with
$x_n+y_n>K$ for all $n\ge N$.

\emph{Step 4 (scratch: what do we need from $x_n$?).} By Step 2, it is
enough that $x_n-B>K$, i.e.\ $x_n>K+B$.

\emph{Step 5 (apply $x_n\to\infty$ with threshold $K+B$).} The definition
of $x_n\to\infty$ holds for every real threshold, in particular for
$M=K+B$. So there is $N$ with $x_n>K+B$ for all $n\ge N$.

\emph{Step 6 (verify).} For $n\ge N$:
$x_n+y_n\ge x_n-B>(K+B)-B=K$.

\emph{Step 7 (conclude).} $K$ was arbitrary, so $x_n+y_n\to\infty$.
\end{proof}
\noindent\key{Only the lower bound $y_n\ge-B$ is used.} So the result holds
whenever $(y_n)$ is merely bounded \emph{below}.
\end{ex}

% =====================================================================
\section{Limit theorems (Week 3--4)}

\begin{thm}
Suppose $x_n\to x$ and $y_n\to y$. Then:
\begin{enumerate}[nosep]
  \item \textbf{Uniqueness:} limits are unique.
  \item \textbf{Bounded:} a convergent sequence is bounded.
  \item $ax_n+by_n\to ax+by$.
  \item $x_ny_n\to xy$.
  \item $x_n/y_n\to x/y$ if $y\ne0$.
  \item \textbf{Squeeze:} if $a_n\le z_n\le b_n$ and $a_n,b_n\to L$, then
        $z_n\to L$.
  \item \textbf{Order:} if $x_n\le y_n$ for all $n$ (large), then
        $x\le y$. In particular, $x_n\ge0\Rightarrow x\ge0$.
        \key{Strict inequality is not preserved:} $\frac1n>0$ but the limit
        is $0$.
\end{enumerate}
\end{thm}

\noindent\key{Proofs to know (short ones are exam favorites):}
\begin{itemize}
  \item \emph{Uniqueness:} $|x-x'|\le|x-x_n|+|x_n-x'|<2\eps$ for every
        $\eps$, so $x=x'$.
  \item \emph{Bounded:} with $\eps=1$, $|x_n|<|x|+1$ for $n\ge N$. Take
        $M=\max\{|x_1|,\dots,|x_{N-1}|,|x|+1\}$.
  \item \emph{Sum:} $|(x_n+y_n)-(x+y)|\le|x_n-x|+|y_n-y|<\frac\eps2+\frac\eps2$,
        using $N=\max(N_1,N_2)$.
  \item \emph{Product:} $|x_ny_n-xy|=|(x_n-x)y_n+x(y_n-y)|\le M|x_n-x|+|x||y_n-y|$,
        where $|y_n|\le M$ because convergent sequences are bounded.
  \item \emph{Squeeze:} $L-\eps<a_n\le z_n\le b_n<L+\eps$.
  \item \emph{Order, by contradiction:} if $x_n\ge0$ and $x<0$, take
        $\eps=-x/2$. Then $x_n<x+\eps=x/2<0$ for large $n$, contradiction.
\end{itemize}

\subsection{The quotient rule and the ``bounded away from zero'' lemma}
\begin{thm}[Lemma]
If $y_n\to y\ne0$, then there is $N_0$ such that $|y_n|>\frac{|y|}{2}$ for
all $n\ge N_0$.
\end{thm}
\begin{proof}
Apply the definition with $\eps=\frac{|y|}2$. Then
$|y_n|\ge|y|-|y-y_n|>|y|-\frac{|y|}2=\frac{|y|}2$ by the reverse triangle
inequality.
\end{proof}
\noindent\emph{Quotient rule sketch:} by the product rule, it suffices to
show $\frac1{y_n}\to\frac1y$. For $n\ge N_0$,
\[
\left|\frac1{y_n}-\frac1y\right|=\frac{|y-y_n|}{|y_n||y|}<\frac{2}{|y|^2}|y_n-y|,
\]
which can be made $<\eps$ by choosing $N$ with $|y_n-y|<\frac{|y|^2\eps}{2}$.

\subsection{More limit tools}
\begin{itemize}
  \item \textbf{Zero times bounded:} if $x_n\to0$ and $|y_n|\le M$, then
        $x_ny_n\to0$, because $|x_ny_n|\le M|x_n|$. (This does \emph{not}
        need $y_n$ to converge.)
  \item \textbf{Shifts don't matter:} $x_n\to L\iff x_{n+1}\to L\iff x_{n+k}\to L$.
        Changing finitely many terms doesn't change the limit.
  \item \textbf{Reciprocals and infinity:} if $x_n\to\infty$, then
        $\frac1{x_n}\to0$. If $x_n>0$ and $x_n\to0$, then $\frac1{x_n}\to\infty$.
  \item \textbf{Infinity arithmetic:} $x_n\to\infty$ and $y_n\ge c>0$
        eventually $\Rightarrow x_ny_n\to\infty$. $x_n\to\infty$ and
        $y_n$ bounded below $\Rightarrow x_n+y_n\to\infty$.
        \key{$\infty-\infty$, $0\cdot\infty$ and $\frac\infty\infty$ are
        indeterminate:} for example $n-n=0$, $(n+1)-n=1$, and $n^2-n\to\infty$.
  \item \textbf{Comparison for infinity:} if $x_n\le y_n$ and
        $x_n\to\infty$, then $y_n\to\infty$.
  \item \textbf{The ``max of the $N$'s'':} when you need two conditions at
        once, take $N=\max(N_1,N_2)$.
  \item \textbf{$\eps$-splitting:} to get a total of $<\eps$ from two pieces,
        make each piece $<\frac\eps2$. With a constant $C$ in front, make the
        piece $<\frac{\eps}{C}$.
\end{itemize}

\subsection{Binomial-theorem estimates (Week 4 style)}
For $x\ge0$ and $n\ge2$:
$(1+x)^n=\sum_{k=0}^n\binom nk x^k\ge 1+nx$ and also
$(1+x)^n\ge\binom n2x^2=\frac{n(n-1)}2x^2$.
Keep only the one term you need.

\begin{ex}[$\frac{n}{2^n}\to0$]
$2^n=(1+1)^n\ge\binom n2=\frac{n(n-1)}{2}$, so for $n\ge2$,
$0<\frac n{2^n}\le\frac{2n}{n(n-1)}=\frac{2}{n-1}\to0$. Squeeze.

\begin{proof}[Full proof, step by step]
\emph{Step 1 (binomial theorem).} For $n\ge2$,
\[
2^n=(1+1)^n=\sum_{k=0}^n\binom nk\ge\binom n2=\frac{n(n-1)}2,
\]
because every term of the sum is $\ge0$ and we kept only the $k=2$ term.

\emph{Step 2 (bound the sequence).} For $n\ge2$, $n-1>0$, so taking
reciprocals (both sides positive) gives $\frac1{2^n}\le\frac{2}{n(n-1)}$.
Multiplying by $n>0$:
\[
0<\frac n{2^n}\le\frac{2n}{n(n-1)}=\frac2{n-1}.
\]

\emph{Step 3 ($\frac2{n-1}\to0$).} Let $\eps>0$ and choose $N\in\N$ with
$N>1+\frac2\eps$ (also $N\ge2$). For $n\ge N$, $n-1\ge N-1>\frac2\eps$, so
$0<\frac2{n-1}<\eps$.

\emph{Step 4 (squeeze).} $0\le\frac n{2^n}\le\frac2{n-1}$ for $n\ge2$,
and both outer sequences tend to $0$. By the Squeeze Theorem,
$\frac n{2^n}\to0$. (Or directly: for $n\ge N$,
$\left|\frac n{2^n}\right|\le\frac2{n-1}<\eps$.)
\end{proof}
\noindent\key{Why $\binom n2$ and not $1+n$?} Bernoulli gives
$2^n\ge1+n$, so $\frac n{2^n}\le\frac n{n+1}$, which does \emph{not} go to
$0$. You need a term of degree $2$ in $n$ to beat the $n$ on top.
\end{ex}

\begin{ex}[$a^{1/n}\to1$ for $a>1$]
Let $d_n=a^{1/n}-1>0$. Then $a=(1+d_n)^n\ge1+nd_n$, so
$0<d_n\le\frac{a-1}{n}\to0$. For $0<a<1$, apply this to $\frac1a$ and take
reciprocals.

\begin{proof}[Full proof, step by step]
\emph{Step 1 (define the error term).} Let $a>1$ and
$d_n:=a^{1/n}-1$. Since $a>1$, $a^{1/n}>1$ (if $a^{1/n}\le1$ then
$a=(a^{1/n})^n\le1$). So $d_n>0$.

\emph{Step 2 (rewrite $a$).} $a^{1/n}=1+d_n$, so $a=(1+d_n)^n$.

\emph{Step 3 (Bernoulli).} Since $d_n>0\ge-1$, Bernoulli's inequality
gives $(1+d_n)^n\ge1+nd_n$. So $a\ge1+nd_n$.

\emph{Step 4 (solve for $d_n$).} $nd_n\le a-1$, so
$0<d_n\le\frac{a-1}n$.

\emph{Step 5 ($\eps$--$N$).} Let $\eps>0$ and choose $N>\frac{a-1}\eps$.
For $n\ge N$:
$|a^{1/n}-1|=d_n\le\frac{a-1}n\le\frac{a-1}N<\eps$. So $a^{1/n}\to1$.

\emph{Step 6 (the case $a=1$).} $1^{1/n}=1$ for all $n$, so the limit is $1$.

\emph{Step 7 (the case $0<a<1$).} Then $b:=\frac1a>1$, so $b^{1/n}\to1$ by
Step 5. Since $a^{1/n}=\frac1{b^{1/n}}$ and the limit $1\ne0$, the quotient
rule gives $a^{1/n}\to\frac11=1$.
\end{proof}
\end{ex}

\begin{ex}[$n^{1/n}\to1$]
Let $d_n=n^{1/n}-1\ge0$. Then $n=(1+d_n)^n\ge\frac{n(n-1)}2d_n^2$, so
$d_n^2\le\frac{2}{n-1}$ and $0\le d_n\le\sqrt{\frac2{n-1}}\to0$.

\begin{proof}[Full proof, step by step]
\emph{Step 1 (error term).} Let $d_n:=n^{1/n}-1$. Since $n\ge1$,
$n^{1/n}\ge1$ (as in Step 1 above), so $d_n\ge0$.

\emph{Step 2 (binomial theorem, keep the $k=2$ term).} For $n\ge2$,
\[
n=(1+d_n)^n=\sum_{k=0}^n\binom nk d_n^k\ge\binom n2d_n^2=\frac{n(n-1)}2d_n^2,
\]
since every term is $\ge0$ (here $d_n\ge0$ is used).

\emph{Step 3 (solve for $d_n$).} Divide by $\frac{n(n-1)}2>0$:
$d_n^2\le\frac{2n}{n(n-1)}=\frac2{n-1}$. Since $d_n\ge0$, taking square
roots gives $0\le d_n\le\sqrt{\frac2{n-1}}$.

\emph{Step 4 (scratch).}
$\sqrt{\frac2{n-1}}<\eps\iff\frac2{n-1}<\eps^2\iff n>1+\frac2{\eps^2}$.

\emph{Step 5 ($\eps$--$N$).} Let $\eps>0$ and choose $N\in\N$ with
$N>1+\frac2{\eps^2}$ (so also $N\ge2$). For $n\ge N$:
\[
|n^{1/n}-1|=d_n\le\sqrt{\frac2{n-1}}\le\sqrt{\frac2{N-1}}<\eps.
\]

\emph{Step 6 (conclude).} So $n^{1/n}\to1$.
\end{proof}
\noindent\key{Why not Bernoulli here?} It gives $n\ge1+nd_n$, i.e.\
$d_n\le\frac{n-1}n$, which does not go to $0$. The $\binom n2$ term is
needed because the base $n$ grows.
\end{ex}

\begin{ex}[$\frac{a^n}{n!}\to0$ for any $a$]
Fix an integer $K\ge2|a|$. For $n>K$,
$\frac{|a|^n}{n!}=\frac{|a|^K}{K!}\cdot\prod_{j=K+1}^n\frac{|a|}{j}\le\frac{|a|^K}{K!}\left(\tfrac12\right)^{n-K}\to0$.

\begin{proof}[Full proof, step by step]
\emph{Step 1 (reduce to $|a|$).} $\left|\frac{a^n}{n!}-0\right|=\frac{|a|^n}{n!}$,
so it is enough to show $\frac{|a|^n}{n!}\to0$. If $a=0$ every term is $0$,
so assume $a\ne0$.

\emph{Step 2 (fix a cutoff $K$).} By the Archimedean property choose an
integer $K\ge1$ with $K\ge2|a|$. Then for every $j\ge K+1$,
$j>2|a|$, so $\frac{|a|}{j}<\frac12$.

\emph{Step 3 (split the product).} For $n>K$, write
$|a|^n=|a|^K\cdot|a|^{n-K}$ and $n!=K!\cdot(K+1)(K+2)\cdots n$. So
\[
\frac{|a|^n}{n!}=\frac{|a|^K}{K!}\cdot\frac{|a|}{K+1}\cdot\frac{|a|}{K+2}\cdots\frac{|a|}{n}.
\]
There are exactly $n-K$ factors after the first.

\emph{Step 4 (bound each tail factor).} By Step 2 each of those $n-K$
factors is $<\frac12$. Let $C:=\frac{|a|^K}{K!}$, a constant (it does not
depend on $n$). Then
\[
0<\frac{|a|^n}{n!}\le C\left(\tfrac12\right)^{n-K}=C\,2^K\left(\tfrac12\right)^n.
\]

\emph{Step 5 (the bound tends to $0$).} $\left(\frac12\right)^n\to0$ (basic
limit, $|\frac12|<1$), so $C2^K\left(\frac12\right)^n\to0$ by the constant
multiple rule.

\emph{Step 6 (squeeze).} $0<\frac{|a|^n}{n!}\le C2^K(\frac12)^n$ for $n>K$.
By the Squeeze Theorem (only large $n$ matter), $\frac{|a|^n}{n!}\to0$,
and therefore $\frac{a^n}{n!}\to0$.
\end{proof}
\noindent\key{Idea:} the first $K$ factors can be large, but there are
only finitely many of them, so they form a fixed constant $C$. After that,
each new factor at least halves the term.
\end{ex}

\subsection{Basic limits to memorize}
\begin{enumerate}[nosep]
  \item $\frac1{n^p}\to0$ for $p>0$.
  \item \textbf{Ratio of polynomials of equal degree:} the limit is the
        ratio of leading coefficients. (Divide the top and bottom by $n^k$.)
        If the top has lower degree, the limit is $0$. If the top has higher
        degree, the limit is $\pm\infty$.
  \item $a^n\to0$ if $|a|<1$. (Proof: write $\frac1{|a|}=1+x$ and use
        Bernoulli: $|a|^n\le\frac{1}{1+nx}\to0$.)
  \item $1+a+\dots+a^n=\frac{1-a^{n+1}}{1-a}\to\frac1{1-a}$ if $|a|<1$.
  \item $\frac{\text{bounded}}{\to\infty}\to0$, e.g.\ $\frac{\sin n}{n}$ and
        $\frac{(-1)^n}{n}$.
  \item If $x_n\to L$ then $|x_n|\to|L|$ (reverse triangle inequality).
\end{enumerate}

% =====================================================================
\section{Monotone sequences and the Monotone Convergence Theorem}

\begin{defn}
$(x_n)$ is \textbf{increasing} if $x_n\le x_{n+1}$ for all $n$,
\textbf{decreasing} if $x_n\ge x_{n+1}$ for all $n$, and
\textbf{monotone} if it is one of the two.
\end{defn}

\begin{method}[Showing monotonicity: three tools]
\begin{enumerate}[nosep]
  \item Show $x_{n+1}-x_n\ge0$ (HW3 P4: common denominator, then check the
        sign).
  \item If $x_n>0$, show $\frac{x_{n+1}}{x_n}\ge1$.
  \item For recursive sequences, use induction: $x_n\le x_{n+1}\Rightarrow x_{n+1}\le x_{n+2}$.
\end{enumerate}
\end{method}

\begin{thm}[Monotone Convergence Theorem]
A bounded monotone sequence converges. If it is increasing and bounded
above, then $x_n\to\sup\{x_n\}$.
\end{thm}
\begin{proof}
Let $L=\sup\{x_n\}$, which exists by completeness. Let $\eps>0$. Since
$L-\eps$ is not an upper bound, there is $N$ with $x_N>L-\eps$. For
$n\ge N$: $L-\eps<x_N\le x_n\le L<L+\eps$, so $|x_n-L|<\eps$.
\end{proof}
\noindent\key{This proof is short and is exactly Theorem~\ref{thm:epssup}(ii)
plus monotonicity. Be able to reproduce it.}

\begin{method}[Recursive sequences: the 4-step routine]
\begin{enumerate}[nosep]
  \item Compute $x_1,x_2,x_3$ to guess the direction and the limit.
  \item Prove a bound by induction.
  \item Prove monotonicity (often by induction, or using the bound).
  \item By MCT, $x_n\to L$. Pass to the limit in the recursion using the
        limit laws, and \emph{also} $x_{n+1}\to L$ (it is the same sequence
        shifted). Solve for $L$ and use the bounds to discard wrong roots.
\end{enumerate}
\end{method}

\begin{ex}
$x_1=1$, $x_{n+1}=\sqrt{2+x_n}$. Then $x_n\to2$.
\begin{proof}
\emph{Bounded:} $0<x_n<2$ by induction. Indeed $x_1=1$, and if
$0<x_n<2$ then $0<x_{n+1}=\sqrt{2+x_n}<\sqrt4=2$.

\emph{Increasing:} $x_{n+1}>x_n\iff 2+x_n>x_n^2\iff(2-x_n)(1+x_n)>0$, which
holds because $0<x_n<2$.

\emph{Limit:} by MCT, $x_n\to L$ with $1\le L\le2$. Square the recursion:
$x_{n+1}^2=2+x_n$. Take limits using the product and sum laws:
$L^2=2+L$, so $(L-2)(L+1)=0$. Since $L\ge1$, $L=2$.
\end{proof}

\begin{proof}[Full proof, step by step]
\emph{Step 1 (compute and guess).} $x_1=1$, $x_2=\sqrt3\approx1.73$,
$x_3=\sqrt{2+\sqrt3}\approx1.93$. The terms increase, apparently toward
$2$. (If $L=\sqrt{2+L}$ then $L=2$.)

\emph{Step 2 (bound by induction: $0<x_n<2$ for all $n$).}
Base: $0<x_1=1<2$. Inductive step: assume $0<x_n<2$. Then
$2<2+x_n<4$, and since $\sqrt{\ }$ is increasing,
$0<\sqrt2<\sqrt{2+x_n}<\sqrt4=2$, i.e.\ $0<x_{n+1}<2$.

\emph{Step 3 (increasing).} Fix $n$. Both $x_n$ and $x_{n+1}$ are
positive, so $x_{n+1}>x_n\iff x_{n+1}^2>x_n^2$. Since
$x_{n+1}^2=2+x_n$:
\[
x_{n+1}>x_n\iff2+x_n-x_n^2>0\iff(2-x_n)(1+x_n)>0.
\]
By Step 2, $2-x_n>0$ and $1+x_n>0$, so the product is positive and
$x_{n+1}>x_n$.

\emph{Step 4 (MCT).} $(x_n)$ is increasing and bounded above by $2$, so by
the Monotone Convergence Theorem it converges, say $x_n\to L$.

\emph{Step 5 (bounds on $L$).} $1=x_1\le x_n<2$ for all $n$, so by the
order property of limits $1\le L\le2$.

\emph{Step 6 (pass to the limit).} The recursion gives
$x_{n+1}^2=2+x_n$ for all $n$. The shifted sequence $(x_{n+1})$ also
converges to $L$, so by the product rule $x_{n+1}^2\to L^2$. By the sum
rule $2+x_n\to2+L$. By uniqueness of limits, $L^2=2+L$.

\emph{Step 7 (solve and discard).} $L^2-L-2=(L-2)(L+1)=0$, so $L=2$ or
$L=-1$. Step 5 rules out $-1$. Hence $x_n\to2$.
\end{proof}
\noindent\key{Order matters:} Step 6 is only valid \emph{after} Step 4.
Without knowing the limit exists, ``$L=\sqrt{2+L}$'' is meaningless.
\end{ex}

\begin{ex}[A decreasing recursive sequence]
$x_1=3$, $x_{n+1}=\sqrt{x_n+2}$. \emph{Bound:} $x_n>2$ by induction
($x_{n+1}>\sqrt4=2$). \emph{Decreasing:}
$x_{n+1}<x_n\iff x_n+2<x_n^2\iff(x_n-2)(x_n+1)>0$, which holds. By MCT,
$x_n\to L\ge2$ with $L^2=L+2$, so $L=2$.

\begin{proof}[Full proof, step by step]
\emph{Step 1 (compute and guess).} $x_1=3$, $x_2=\sqrt5\approx2.24$,
$x_3=\sqrt{2+\sqrt5}\approx2.06$. The terms decrease, apparently toward $2$.

\emph{Step 2 (bound by induction: $x_n>2$ for all $n$).} Base: $x_1=3>2$.
Inductive step: assume $x_n>2$. Then $x_n+2>4$, so
$x_{n+1}=\sqrt{x_n+2}>\sqrt4=2$.

\emph{Step 3 (decreasing).} Both $x_n,x_{n+1}>0$, so
$x_{n+1}<x_n\iff x_{n+1}^2<x_n^2\iff x_n+2<x_n^2$, i.e.\
\[
x_n^2-x_n-2=(x_n-2)(x_n+1)>0.
\]
By Step 2, $x_n-2>0$ and $x_n+1>0$, so this holds. Thus $x_{n+1}<x_n$.

\emph{Step 4 (MCT).} $(x_n)$ is decreasing and bounded below by $2$. By
the Monotone Convergence Theorem (decreasing version: it converges to
$\inf\{x_n\}$), $x_n\to L$ for some $L$.

\emph{Step 5 (bounds on $L$).} $2<x_n\le3$ for all $n$, so by the order
property $2\le L\le3$. (Strict $>$ becomes $\ge$ in the limit.)

\emph{Step 6 (pass to the limit).} $x_{n+1}^2=x_n+2$. The left side tends
to $L^2$ (shift and product rule). The right side tends to $L+2$. So
$L^2=L+2$.

\emph{Step 7 (solve and discard).} $(L-2)(L+1)=0$, so $L\in\{2,-1\}$, and
Step 5 forces $L=2$.
\end{proof}
\noindent\key{Compare with the previous example.} Same recursion, but a
different starting point. Starting above the fixed point $2$ gives a
decreasing sequence, and starting below it gives an increasing one. The
bound you prove by induction always points toward the fixed point.
\end{ex}

\begin{ex}[$e$ exists]
$a_n=\left(1+\frac1n\right)^n$ is increasing and bounded above by $3$
(binomial expansion compared with $\sum\frac1{k!}\le1+\sum\frac1{2^{k-1}}$).
So it converges, and its limit is called $e$. \key{MCT proves existence
without knowing the value.}

\begin{proof}[Full proof, step by step]
\emph{Step 1 (expand with the binomial theorem).}
\[
a_n=\sum_{k=0}^n\binom nk\frac1{n^k}.
\]
For $k\ge1$,
\[
\binom nk\frac1{n^k}=\frac{n(n-1)\cdots(n-k+1)}{k!\,n^k}
=\frac1{k!}\cdot\frac nn\cdot\frac{n-1}n\cdots\frac{n-k+1}n
=\frac1{k!}\prod_{j=0}^{k-1}\Bigl(1-\frac jn\Bigr).
\]
So
\[
a_n=\sum_{k=0}^n\frac1{k!}\prod_{j=0}^{k-1}\Bigl(1-\frac jn\Bigr),
\]
where the empty product (for $k=0$) is $1$.

\emph{Step 2 (compare $a_n$ with $a_{n+1}$ term by term).} Fix $k\le n$.
For each $j\ge0$, $\frac j{n+1}\le\frac jn$, so
$1-\frac j{n+1}\ge1-\frac jn\ge0$ (as $j\le k-1<n$). All factors are
nonnegative, so the products compare the same way:
\[
\frac1{k!}\prod_{j=0}^{k-1}\Bigl(1-\frac j{n+1}\Bigr)
\ge\frac1{k!}\prod_{j=0}^{k-1}\Bigl(1-\frac jn\Bigr).
\]
So term $k$ of $a_{n+1}$ is $\ge$ term $k$ of $a_n$, for each
$k=0,\dots,n$.

\emph{Step 3 (increasing).} $a_{n+1}$ also has one extra term ($k=n+1$),
which is $\ge0$. Adding up Step 2 over $k=0,\dots,n$ gives
$a_{n+1}\ge a_n$.

\emph{Step 4 (drop the products).} Each factor
$1-\frac jn\in[0,1]$, so each product is $\le1$ and
$a_n\le\sum_{k=0}^n\frac1{k!}$.

\emph{Step 5 ($k!\ge2^{k-1}$ for $k\ge1$).} By induction: $1!=1=2^0$, and
if $k!\ge2^{k-1}$ then $(k+1)!=(k+1)k!\ge2\cdot2^{k-1}=2^k$ (since
$k+1\ge2$).

\emph{Step 6 (geometric bound).}
\[
a_n\le1+\sum_{k=1}^n\frac1{k!}\le1+\sum_{k=1}^n\frac1{2^{k-1}}
=1+\frac{1-(1/2)^n}{1-1/2}=1+2\Bigl(1-\frac1{2^n}\Bigr)<3.
\]

\emph{Step 7 (MCT).} $(a_n)$ is increasing (Step 3) and bounded above by
$3$ (Step 6). By the Monotone Convergence Theorem it converges. Its limit
is \emph{defined} to be $e$, and $2=a_1\le e\le3$.
\end{proof}
\end{ex}

\begin{ex}[HW3 P4, averages]
If $(x_n)$ is increasing, then so is $y_n=\frac{x_1+\dots+x_n}{n}$. The
numerator of $y_{n+1}-y_n$ is $nx_{n+1}-S_n\ge0$, because each
$x_k\le x_{n+1}$. (Positivity is not actually needed.)

\begin{proof}[Full proof, step by step]
Let $S_n=x_1+\dots+x_n$, so $y_n=\frac{S_n}n$ and $S_{n+1}=S_n+x_{n+1}$.

\emph{Step 1 (goal).} Show $y_{n+1}-y_n\ge0$ for every $n$.

\emph{Step 2 (common denominator).}
\[
y_{n+1}-y_n=\frac{S_{n+1}}{n+1}-\frac{S_n}{n}
=\frac{nS_{n+1}-(n+1)S_n}{n(n+1)}.
\]

\emph{Step 3 (simplify the numerator).} Substitute
$S_{n+1}=S_n+x_{n+1}$:
\[
nS_{n+1}-(n+1)S_n=nS_n+nx_{n+1}-nS_n-S_n=nx_{n+1}-S_n.
\]

\emph{Step 4 (rewrite as a sum of differences).} $nx_{n+1}$ is $x_{n+1}$
added $n$ times, and $S_n$ has $n$ terms, so
\[
nx_{n+1}-S_n=\sum_{k=1}^n\bigl(x_{n+1}-x_k\bigr).
\]

\emph{Step 5 (each difference is $\ge0$).} Since $(x_n)$ is increasing,
$x_k\le x_{k+1}\le\dots\le x_{n+1}$ for each $k\le n$. (Formally, by
chaining $x_k\le x_{k+1}$ repeatedly, or by induction on $n+1-k$.) So
every term $x_{n+1}-x_k\ge0$, and the sum is $\ge0$.

\emph{Step 6 (sign of the fraction).} The numerator is $\ge0$ (Step 5) and
the denominator $n(n+1)>0$. So $y_{n+1}-y_n\ge0$.

\emph{Step 7 (conclude).} $y_n\le y_{n+1}$ for all $n$, so $(y_n)$ is
increasing.
\end{proof}
\noindent\key{Intuition:} adding a new value $x_{n+1}$ that is at least as
large as every earlier value cannot lower the average.
\end{ex}

% =====================================================================
\section{Subsequences and Cauchy sequences (Week 5)}
\label{sec:cauchy}

\subsection{Subsequence facts you need}
\begin{itemize}[nosep]
  \item A \textbf{subsequence} is $(x_{n_k})_k$ with $n_1<n_2<\cdots$.
        Its index is $k$, and \key{$n_k\ge k$ for every $k$}
        (Lemma~\ref{lem:nkk}).
  \item If $x_n\to L$ then every subsequence $\to L$. \key{So two
        subsequences with different limits imply divergence}, e.g.\
        $(-1)^n$ with its even and odd terms.
  \item Every sequence has a monotone subsequence
        (Lemma~\ref{lem:monosub}).
  \item \textbf{Bolzano--Weierstrass:} every bounded sequence has a
        convergent subsequence (Lemma~\ref{lem:bw}).
\end{itemize}

\subsection{Definition and intuition}

\begin{defn}[Cauchy sequence]
$(x_n)$ is \textbf{Cauchy} if for every $\eps>0$ there is $N\in\N$ such that
\[
|x_n-x_m|<\eps\qquad\text{for all } n,m\ge N.
\]
(The lecture notes write $n,m>N$. The two versions are equivalent: replace
$N$ by $N+1$.)
\end{defn}

\noindent\key{Lecture wording (Week 5, Definition 6) --- memorize this.}
\begin{center}
\fbox{\parbox{0.9\textwidth}{A sequence $(x_n)_n$ of real numbers is called
\textbf{Cauchy} if for any $\eps>0$ there exists $N$ so that
$|x_n-x_m|<\eps$ for all $n,m>N$.}}
\end{center}

\noindent\key{In symbols.}
\[
(x_n)\text{ is Cauchy}\iff
\forall\eps>0\ \ \exists N\in\N\ \ \forall n,m\ge N:\ \ |x_n-x_m|<\eps.
\]

\noindent\key{Reading the quantifiers in order.}
\begin{itemize}[nosep]
  \item $\forall\eps>0$: the challenger names any tolerance, however
        small.
  \item $\exists N$: you answer with a cutoff. \textbf{$N$ may depend on
        $\eps$}, but not on $n$ or $m$.
  \item $\forall n,m\ge N$: \emph{every} pair of terms past the cutoff,
        not just neighbors and not just $m=n+1$.
  \item $|x_n-x_m|<\eps$: those two terms are within $\eps$ of each other.
\end{itemize}

\noindent\key{Side by side with convergence.}
\begin{center}
\renewcommand{\arraystretch}{1.3}
\begin{tabular}{|l|l|}
\hline
$x_n\to L$ & $\forall\eps>0\ \exists N\ \forall n\ge N:\ |x_n-L|<\eps$\\ \hline
$(x_n)$ Cauchy & $\forall\eps>0\ \exists N\ \forall n,m\ge N:\ |x_n-x_m|<\eps$\\ \hline
\end{tabular}
\end{center}
The only change: the fixed target $L$ is replaced by a second term
$x_m$. That is why the Cauchy condition can be checked without knowing
$L$.

\noindent\key{Equivalent form.} Writing $m=n+p$: $(x_n)$ is Cauchy iff
$\forall\eps>0\ \exists N$ such that $|x_{n+p}-x_n|<\eps$ for all
$n\ge N$ and \emph{all} $p\in\N$. The ``all $p$'' is essential. Fixing
$p=1$ gives the weaker condition in the Warning below.

\noindent\key{Proof skeleton (copy this shape).}
\begin{quote}
Let $\eps>0$. Choose $N=\ldots$ \emph{(from scratch work)}. Let
$n,m\ge N$; without loss of generality $m>n$. Then
$|x_m-x_n|\le\ldots\le\ldots<\eps$. Hence $(x_n)$ is Cauchy.
\end{quote}
Example~\ref{ex:cdirect} below fills this in for $x_n=\frac n{n+1}$.

\noindent\key{Intuition.} Convergent means ``the terms get close to
\emph{a number $L$}.'' Cauchy means ``the terms get close to \emph{each
other}.'' \textbf{No $L$ appears in the definition.} That is the whole
point: you can prove a sequence converges without knowing or guessing its
limit, even if it is not monotone (so MCT does not apply).

\medskip
\noindent\key{Negation (you need this to prove ``not Cauchy'').}
\[
(x_n)\text{ is not Cauchy}\iff
\exists\eps>0\ \forall N\ \exists n,m\ge N:\ |x_n-x_m|\ge\eps.
\]
In words: some fixed gap $\eps$ keeps appearing between pairs of terms
arbitrarily far out.

\medskip
\noindent\key{Warning.} Cauchy is \emph{not} the same as
$|x_{n+1}-x_n|\to0$. The definition controls \emph{all} pairs $n,m\ge N$,
not only neighbors. See Example~\ref{ex:sqrtn}.

\subsection{Methods}

\begin{method}[Proving Cauchy directly]\label{meth:cauchy}
\leavevmode
\begin{enumerate}[nosep]
  \item By symmetry assume $m>n$ (if $m=n$ the difference is $0$).
  \item \textbf{Scratch work:} simplify or estimate $|x_m-x_n|$ by an
        expression that \key{depends only on $n$} (not on $m$) and goes to
        $0$, e.g.\ $\frac Cn$ or $Cr^n$.
  \item Choose $N$ so that this expression is $<\eps$ when $n\ge N$.
  \item Write the forward proof: ``Let $\eps>0$. Choose $N=\ldots$. Let
        $m>n\ge N$. Then $|x_m-x_n|\le\ldots<\eps$.''
\end{enumerate}
\end{method}

\begin{method}[Telescoping and geometric sums]\label{meth:geom}
If $|x_{k+1}-x_k|\le Cr^k$ for all $k$, with $0<r<1$, then for $m>n$
\[
|x_m-x_n|
=\Bigl|\sum_{k=n}^{m-1}(x_{k+1}-x_k)\Bigr|
\le\sum_{k=n}^{m-1}|x_{k+1}-x_k|
\le C\sum_{k=n}^{m-1}r^k
< C\,\frac{r^n}{1-r}.
\]
Since $r^n\to0$, this is $<\eps$ for large $n$, so $(x_n)$ is Cauchy. Use
this for recursive sequences that oscillate, where MCT fails
(Example~\ref{ex:golden}).
\end{method}

\begin{method}[Proving not Cauchy]
Find one $\eps$ and, for each $N$, a specific pair $n,m\ge N$ with
$|x_n-x_m|\ge\eps$. Good pairs to try: $m=n+1$ (oscillating sequences),
$m=2n$ (sums and slowly growing sequences).
\end{method}

\begin{method}[Which tool?]
\begin{itemize}[nosep]
  \item You know or can guess $L$: $\eps$--$N$.
  \item You can't guess $L$, but the sequence is monotone and bounded: MCT.
  \item You can't guess $L$ and the sequence is not monotone: \textbf{show
        it is Cauchy}.
  \item To show divergence: two subsequences with different limits, or
        show it is not Cauchy, or show it is unbounded.
\end{itemize}
\end{method}

\subsection{Worked examples}

\begin{ex}[A direct Cauchy proof]\label{ex:cdirect}
Prove that $x_n=\dfrac{n}{n+1}$ is Cauchy.

\begin{scratch}
$x_n=1-\frac1{n+1}$. For $m>n$:
$|x_m-x_n|=\bigl|\frac1{n+1}-\frac1{m+1}\bigr|=\frac1{n+1}-\frac1{m+1}<\frac1{n+1}<\frac1n$.
This depends only on $n$. Need $\frac1n<\eps$, i.e.\ $n>\frac1\eps$.
\end{scratch}

\begin{proof}
Let $\eps>0$. By the Archimedean property choose $N\in\N$ with
$N>\frac1\eps$. Let $n,m\ge N$; by symmetry assume $m>n$ (if $m=n$ there
is nothing to prove). Since $x_k=1-\frac1{k+1}$,
\[
|x_m-x_n|=\frac1{n+1}-\frac1{m+1}<\frac1{n+1}<\frac1n\le\frac1N<\eps.
\]
So $(x_n)$ is Cauchy.
\end{proof}
\end{ex}

\begin{ex}[Partial sums of $\sum\frac1{k^2}$ are Cauchy]\label{ex:sumsq}
Let $s_n=\sum_{k=1}^n\frac1{k^2}$. Prove $(s_n)$ is Cauchy.

\begin{scratch}
For $m>n$, $s_m-s_n=\sum_{k=n+1}^m\frac1{k^2}$. We can't sum this exactly,
but $\frac1{k^2}\le\frac1{k(k-1)}=\frac1{k-1}-\frac1k$ for $k\ge2$, which
telescopes.
\end{scratch}

\begin{proof}
Let $m>n\ge1$. Every $k$ in the sum satisfies $k\ge n+1\ge2$, so
$k^2\ge k(k-1)>0$ and
\[
0<s_m-s_n=\sum_{k=n+1}^m\frac1{k^2}\le\sum_{k=n+1}^m\Bigl(\frac1{k-1}-\frac1k\Bigr)
=\frac1n-\frac1m<\frac1n.
\]
Let $\eps>0$ and choose $N>\frac1\eps$. For $m>n\ge N$:
$|s_m-s_n|<\frac1n\le\frac1N<\eps$. So $(s_n)$ is Cauchy, and therefore
converges, even though we never found its limit. (It is
$\frac{\pi^2}{6}$, which you would not be expected to know.)
\end{proof}
\noindent Here MCT would also work ($s_n$ is increasing and $s_n<2$), but
the Cauchy proof doesn't need monotonicity, so it also handles
$\sum\frac{(-1)^k}{k^2}$ or $\sum\frac{\cos k}{k^2}$ with the same estimate.
\end{ex}

\begin{ex}[The harmonic sum is not Cauchy]\label{ex:harm}
Let $h_n=\sum_{k=1}^n\frac1k$. Prove $(h_n)$ is not Cauchy, so it diverges.

\begin{proof}
Take $\eps=\frac12$. Let $N\in\N$ be arbitrary, and set $n=N$, $m=2N$, so
$n,m\ge N$. The sum $h_{2N}-h_N=\sum_{k=N+1}^{2N}\frac1k$ has $N$ terms,
each $\ge\frac1{2N}$ (because $k\le2N$). So
\[
|h_{2N}-h_N|\ge N\cdot\frac1{2N}=\frac12=\eps.
\]
This is the negation of Cauchy. So $(h_n)$ is not Cauchy and hence not
convergent. (Since it is increasing and not convergent, MCT forces it to
be unbounded, so in fact $h_n\to\infty$.)
\end{proof}
\end{ex}

\begin{ex}[Neighbors close $\not\Rightarrow$ Cauchy]\label{ex:sqrtn}
Let $x_n=\sqrt n$. Then
$x_{n+1}-x_n=\frac1{\sqrt{n+1}+\sqrt n}\to0$, but $(x_n)$ is not Cauchy.

\begin{proof}
Take $\eps=1$. Given $N$, let $n=N$ and $m=4N$. Then
$|x_m-x_n|=2\sqrt N-\sqrt N=\sqrt N\ge1=\eps$.
\end{proof}
\noindent \key{Lesson:} consecutive differences can shrink while the
terms still drift off to $\infty$. On a true/false question,
``$x_{n+1}-x_n\to0\Rightarrow$ Cauchy'' is \textbf{false}.
\end{ex}

\begin{ex}[$(-1)^n$ is not Cauchy]
Take $\eps=2$ (or anything $\le2$). For any $N$, $n=N$ and $m=N+1$ give
$|x_m-x_n|=|(-1)^{N+1}-(-1)^N|=2$. So $(-1)^n$ is bounded but not
Cauchy, which shows that \textbf{bounded $\not\Rightarrow$ Cauchy}.
\end{ex}

\begin{ex}[An oscillating recursive sequence: Cauchy beats MCT]\label{ex:golden}
Let $x_1=1$ and $x_{n+1}=1+\dfrac1{x_n}$. Prove $(x_n)$ converges and find
the limit.

\smallskip\noindent\textbf{Why MCT fails.} The terms are
$1,\ 2,\ \frac32,\ \frac53,\ \frac85,\ \dots$, which go up, down, up,
down. The sequence is not monotone.

\begin{proof}
\emph{Step 1 ($x_n\ge1$ for all $n$).} Induction: $x_1=1$. If $x_n\ge1$,
then $\frac1{x_n}>0$, so $x_{n+1}=1+\frac1{x_n}>1$.

\emph{Step 2 (products of neighbors are $\ge2$).} For $n\ge2$,
$x_nx_{n-1}=\bigl(1+\frac1{x_{n-1}}\bigr)x_{n-1}=x_{n-1}+1\ge2$ by Step 1.

\emph{Step 3 (contraction).} For $n\ge2$,
\[
|x_{n+1}-x_n|=\Bigl|\frac1{x_n}-\frac1{x_{n-1}}\Bigr|
=\frac{|x_{n-1}-x_n|}{x_nx_{n-1}}\le\frac12|x_n-x_{n-1}|.
\]

\emph{Step 4 (geometric bound).} Since $|x_2-x_1|=1$, induction on Step 3
gives $|x_{k+1}-x_k|\le\bigl(\frac12\bigr)^{k-1}$ for all $k\ge1$.

\emph{Step 5 (Cauchy).} For $m>n$, by Method~\ref{meth:geom},
\[
|x_m-x_n|\le\sum_{k=n}^{m-1}\Bigl(\frac12\Bigr)^{k-1}
=\Bigl(\frac12\Bigr)^{n-1}\cdot\frac{1-(1/2)^{m-n}}{1/2}
<\Bigl(\frac12\Bigr)^{n-2}=\frac4{2^n}\le\frac4n
\]
(using $2^n\ge n$). Given $\eps>0$, choose $N>\frac4\eps$. Then
$|x_m-x_n|<\eps$ for $m>n\ge N$. So $(x_n)$ is Cauchy.

\emph{Step 6 (find the limit).} By the Cauchy criterion, $x_n\to L$ for
some $L$, and $L\ge1$ by Step 1 (limits preserve $\ge$). In particular
$L\ne0$, so the quotient rule applies to $x_{n+1}=1+\frac1{x_n}$. Taking
limits on both sides (note $x_{n+1}\to L$ too) gives $L=1+\frac1L$, i.e.\
$L^2-L-1=0$, so $L=\frac{1\pm\sqrt5}2$. Since $L\ge1$,
$L=\frac{1+\sqrt5}2$.
\end{proof}
\end{ex}

\begin{ex}[Sums and products of Cauchy sequences]\label{ex:sumprod}
Let $(x_n)$ and $(y_n)$ be Cauchy. Prove directly from the definition
(without the Cauchy criterion) that $(x_n+y_n)$ and $(x_ny_n)$ are Cauchy.

\begin{proof}
\emph{Sum.} Let $\eps>0$. Choose $N_1$ with $|x_n-x_m|<\frac\eps2$ for
$n,m\ge N_1$, and $N_2$ with $|y_n-y_m|<\frac\eps2$ for $n,m\ge N_2$. Let
$N=\max\{N_1,N_2\}$. For $n,m\ge N$,
\[
|(x_n+y_n)-(x_m+y_m)|\le|x_n-x_m|+|y_n-y_m|<\tfrac\eps2+\tfrac\eps2=\eps.
\]

\emph{Product.} Cauchy sequences are bounded (Lemma~\ref{lem:cbdd}), so
there is $M>0$ with $|x_n|\le M$ and $|y_n|\le M$ for all $n$. Add and
subtract $x_ny_m$:
\[
|x_ny_n-x_my_m|=|x_n(y_n-y_m)+y_m(x_n-x_m)|\le M|y_n-y_m|+M|x_n-x_m|.
\]
Let $\eps>0$ and choose $N$ so that $|x_n-x_m|<\frac\eps{2M}$ and
$|y_n-y_m|<\frac\eps{2M}$ for $n,m\ge N$. Then
$|x_ny_n-x_my_m|<\frac\eps2+\frac\eps2=\eps$.
\end{proof}
\end{ex}

\begin{ex}[Why completeness matters: Cauchy in $\Q$]
Let $q_1=1$, $q_2=1.4$, $q_3=1.41$, $q_4=1.414,\dots$ be the decimal
truncations of $\sqrt2$. Each $q_n\in\Q$ and $|q_m-q_n|<10^{-(n-1)}$ for
$m>n$, so $(q_n)$ is Cauchy. Its limit in $\R$ is $\sqrt2\notin\Q$. So
\textbf{in $\Q$ a Cauchy sequence need not converge}. The proof of the
Cauchy criterion below has to use the Completeness Axiom somewhere. It
enters through MCT, which needs the existence of $\sup$.
\end{ex}

% ---------------------------------------------------------------------
\subsection{The long proof: the Cauchy criterion, every step}

\begin{thm}[Cauchy criterion]\label{thm:cauchy}
A sequence $(x_n)$ of real numbers converges if and only if it is Cauchy.
\end{thm}

\noindent\key{Proof map.} ($\Rightarrow$) is a short triangle-inequality
argument. ($\Leftarrow$) is built from a chain of lemmas:
\begin{align*}
\text{Cauchy}
&\xrightarrow{\ \text{Lem.~\ref{lem:cbdd}}\ }\text{bounded}
\xrightarrow{\ \text{Lem.~\ref{lem:monosub}}\ }\text{bounded monotone subsequence}\\
&\xrightarrow{\ \text{MCT}\ }\text{convergent subsequence}
\xrightarrow{\ \text{Lem.~\ref{lem:csub}}\ }\text{convergent}.
\end{align*}
Lemmas~\ref{lem:monosub} and~\ref{lem:bw} together are
Bolzano--Weierstrass. Each lemma below is proved in full.

\subsubsection*{Part I: convergent $\Rightarrow$ Cauchy}

\begin{proof}[Proof of ($\Rightarrow$)]
\emph{Step 1 (setup).} Suppose $x_n\to L$. We must show: for every
$\eps>0$ there is $N$ such that $|x_n-x_m|<\eps$ for all $n,m\ge N$.

\emph{Step 2 (scratch: how to compare two terms).} We know nothing directly
about $x_n-x_m$, but we know both terms are close to $L$. Insert $L$:
$x_n-x_m=(x_n-L)+(L-x_m)$. Each piece will get half of the budget $\eps$.

\emph{Step 3 (use convergence with $\frac\eps2$).} Let $\eps>0$. Then
$\frac\eps2>0$, so by the definition of $x_n\to L$ there is $N\in\N$ with
\[
|x_k-L|<\frac\eps2\qquad\text{for all }k\ge N.
\]

\emph{Step 4 (apply it twice).} Let $n,m\ge N$. Applying Step 3 with $k=n$
and with $k=m$ gives $|x_n-L|<\frac\eps2$ and $|x_m-L|<\frac\eps2$.

\emph{Step 5 (triangle inequality).}
\[
|x_n-x_m|=|(x_n-L)+(L-x_m)|\le|x_n-L|+|L-x_m|<\frac\eps2+\frac\eps2=\eps.
\]

\emph{Step 6 (conclude).} $\eps>0$ was arbitrary, and the $N$ from Step 3
works for all $n,m\ge N$. So $(x_n)$ is Cauchy.
\end{proof}

\subsubsection*{Part II: Cauchy $\Rightarrow$ convergent}

\begin{lem}[Cauchy sequences are bounded]\label{lem:cbdd}
If $(x_n)$ is Cauchy, there is $M>0$ with $|x_n|\le M$ for all $n\in\N$.
\end{lem}

\begin{proof}
\emph{Step 1 (pick a specific $\eps$).} Boundedness doesn't involve any
$\eps$, so we may use the Cauchy property with any fixed tolerance we
like. Take $\eps=1$: there is $N\in\N$ with $|x_n-x_m|<1$ for all
$n,m\ge N$.

\emph{Step 2 (anchor at $x_N$).} Fix $m=N$. For every $n\ge N$ we get
$|x_n-x_N|<1$.

\emph{Step 3 (bound the tail).} For $n\ge N$, by the triangle inequality,
\[
|x_n|=|(x_n-x_N)+x_N|\le|x_n-x_N|+|x_N|<1+|x_N|.
\]

\emph{Step 4 (handle the finitely many early terms).} The terms
$x_1,\dots,x_{N-1}$ are finitely many real numbers, so they have a largest
absolute value. Define
\[
M=\max\bigl\{|x_1|,\ |x_2|,\ \dots,\ |x_{N-1}|,\ 1+|x_N|\bigr\}.
\]
(If $N=1$ the list is just $1+|x_1|$.) Then $M\ge1>0$.

\emph{Step 5 (check every $n$).} If $n<N$, then $|x_n|$ is one of the
numbers in the max, so $|x_n|\le M$. If $n\ge N$, then
$|x_n|<1+|x_N|\le M$ by Step 3. So $|x_n|\le M$ for all $n$.
\end{proof}
\noindent\key{Why the max?} The Cauchy condition only controls the tail
$n\ge N$. A finite list can always be bounded by its maximum, but an
infinite list can't. That's why the proof splits into ``early'' and
``tail'' terms. The same trick proves ``convergent $\Rightarrow$
bounded.''

\begin{lem}[Indices of a subsequence]\label{lem:nkk}
If $n_1<n_2<\cdots$ are natural numbers, then $n_k\ge k$ for all
$k\in\N$.
\end{lem}

\begin{proof}
Induction on $k$. \emph{Base:} $n_1\in\N$, so $n_1\ge1$.
\emph{Step:} assume $n_k\ge k$. Then $n_{k+1}>n_k\ge k$, so
$n_{k+1}>k$. Since $n_{k+1}$ and $k$ are integers, $n_{k+1}\ge k+1$.
\end{proof}

\begin{lem}[Monotone Subsequence Theorem]\label{lem:monosub}
Every sequence $(x_n)$ of real numbers has a monotone subsequence.
\end{lem}

\begin{proof}
\emph{Step 1 (key definition).} Call an index $p\in\N$ a \textbf{peak} if
$x_p\ge x_n$ for every $n>p$. Think of standing at position $p$ and
looking right: nothing ahead is higher than $x_p$.

\emph{Step 2 (two exhaustive cases).} Either there are infinitely many
peaks, or there are finitely many (possibly none). We find a monotone
subsequence in each case.

\emph{Step 3 (Case 1: infinitely many peaks).} List the peaks in
increasing order, $p_1<p_2<p_3<\cdots$. (Any infinite subset of $\N$ can
be listed this way by repeatedly taking the smallest remaining element,
using well-ordering.) Set $n_k=p_k$.
\begin{itemize}[nosep]
  \item \emph{Subsequence:} $n_1<n_2<\cdots$ by construction.
  \item \emph{Monotone:} fix $k$. Since $p_k$ is a peak, $x_{p_k}\ge x_n$
        for every $n>p_k$. In particular $p_{k+1}>p_k$, so
        $x_{p_k}\ge x_{p_{k+1}}$.
\end{itemize}
So $(x_{n_k})$ is \textbf{decreasing}.

\emph{Step 4 (Case 2: finitely many peaks; setup).} Choose $N_0\in\N$
larger than every peak (if there are no peaks, take $N_0=1$). Then
\key{no index $n\ge N_0$ is a peak}.

\emph{Step 5 (negate ``peak'').} ``$p$ is a peak'' means
$\forall n>p:\ x_p\ge x_n$. So ``$p$ is \emph{not} a peak'' means
\[
\exists n>p\ \text{ with }\ x_n>x_p.
\]

\emph{Step 6 (recursive construction).} Let $n_1=N_0$. Suppose
$n_k\ge N_0$ has been chosen. Then $n_k$ is not a peak, so by Step 5 the
set $\{n>n_k: x_n>x_{n_k}\}$ is nonempty. Let $n_{k+1}$ be its smallest
element (well-ordering makes the choice definite). Then
\[
n_{k+1}>n_k\qquad\text{and}\qquad x_{n_{k+1}}>x_{n_k}.
\]
Also $n_{k+1}>n_k\ge N_0$, so $n_{k+1}$ is again not a peak, and the
construction continues. By induction, $n_k$ is defined for every $k$.

\emph{Step 7 (check).} The indices are strictly increasing, so
$(x_{n_k})$ is a subsequence, and $x_{n_1}<x_{n_2}<\cdots$, so it is
(strictly) \textbf{increasing}.

\emph{Step 8 (conclude).} In both cases $(x_n)$ has a monotone
subsequence.
\end{proof}

\begin{lem}[Bolzano--Weierstrass]\label{lem:bw}
Every bounded sequence has a convergent subsequence.
\end{lem}

\begin{proof}
\emph{Step 1.} Let $|x_n|\le M$ for all $n$. By
Lemma~\ref{lem:monosub}, $(x_n)$ has a monotone subsequence $(x_{n_k})$.
(Boundedness has not been used yet.)

\emph{Step 2 (the subsequence is bounded).} Each $x_{n_k}$ is a term of
the original sequence, so $|x_{n_k}|\le M$ for all $k$.

\emph{Step 3 (MCT).} $(x_{n_k})$ is monotone and bounded, so by the
Monotone Convergence Theorem it converges. \key{This is where
completeness enters}: MCT says the limit is the sup (or inf) of the
terms, and that sup exists by the Completeness Axiom.
\end{proof}

\begin{lem}[A Cauchy sequence follows its convergent subsequence]\label{lem:csub}
If $(x_n)$ is Cauchy and some subsequence $x_{n_k}\to L$, then
$x_n\to L$.
\end{lem}

\begin{proof}
\emph{Step 1 (goal).} Let $\eps>0$. We need $N$ with $|x_n-L|<\eps$ for
all $n\ge N$.

\emph{Step 2 (scratch: insert a subsequence term).} We control
$|x_n-x_m|$ (Cauchy) and $|x_{n_k}-L|$ (subsequence). Insert $x_{n_k}$:
$|x_n-L|\le|x_n-x_{n_k}|+|x_{n_k}-L|$, and give each piece
$\frac\eps2$.

\emph{Step 3 (Cauchy with $\frac\eps2$).} There is $N_1$ with
$|x_n-x_m|<\frac\eps2$ for all $n,m\ge N_1$.

\emph{Step 4 (subsequence with $\frac\eps2$).} There is $K$ with
$|x_{n_k}-L|<\frac\eps2$ for all $k\ge K$.

\emph{Step 5 (pick one good subsequence index).} Let
$k_0=\max\{K,N_1\}$. Then $k_0\ge K$, so $|x_{n_{k_0}}-L|<\frac\eps2$.
By Lemma~\ref{lem:nkk}, $n_{k_0}\ge k_0\ge N_1$. Note that $k_0$ depends
only on $\eps$, not on $n$.

\emph{Step 6 (choose $N$ and verify).} Let $N=N_1$ and let $n\ge N$.
Both $n$ and $n_{k_0}$ are $\ge N_1$, so Step 3 with $m=n_{k_0}$ gives
$|x_n-x_{n_{k_0}}|<\frac\eps2$. Therefore
\[
|x_n-L|\le|x_n-x_{n_{k_0}}|+|x_{n_{k_0}}-L|<\frac\eps2+\frac\eps2=\eps.
\]

\emph{Step 7.} $\eps$ was arbitrary, so $x_n\to L$.
\end{proof}

\begin{proof}[Proof of Theorem~\ref{thm:cauchy}, ($\Leftarrow$)]
Let $(x_n)$ be Cauchy.
\begin{enumerate}[nosep]
  \item By Lemma~\ref{lem:cbdd}, $(x_n)$ is bounded.
  \item By Bolzano--Weierstrass (Lemma~\ref{lem:bw}), it has a convergent
        subsequence $x_{n_k}\to L$ for some $L\in\R$.
  \item By Lemma~\ref{lem:csub}, $x_n\to L$.
\end{enumerate}
Together with Part I, this proves Theorem~\ref{thm:cauchy}.
\end{proof}

\subsubsection*{Exam-length version (what to write in about 10 minutes)}
If you are asked to ``prove every Cauchy sequence converges'' and may
quote B--W:

\begin{proof}
Let $(x_n)$ be Cauchy. With $\eps=1$ there is $N$ such that $|x_n-x_N|<1$
for $n\ge N$, so $|x_n|\le\max\{|x_1|,\dots,|x_{N-1}|,1+|x_N|\}$ for all
$n$, and $(x_n)$ is bounded. By Bolzano--Weierstrass some subsequence
$x_{n_k}\to L$. Let $\eps>0$. Choose $N_1$ with $|x_n-x_m|<\frac\eps2$ for
$n,m\ge N_1$, and $K$ with $|x_{n_k}-L|<\frac\eps2$ for $k\ge K$. Fix
$k=\max\{K,N_1\}$, so $n_k\ge k\ge N_1$. For $n\ge N_1$,
$|x_n-L|\le|x_n-x_{n_k}|+|x_{n_k}-L|<\eps$. So $x_n\to L$.
\end{proof}

\noindent\key{Points graders look for:} (1) boundedness proved, with the
finite-max trick; (2) B--W named; (3) the $\frac\eps2$ split; (4) the
justification $n_k\ge k\ge N_1$.

% ---------------------------------------------------------------------
\subsection{The Nested Interval Theorem}
\label{sec:nit}

\begin{thm}[Nested Interval Theorem]\label{thm:nit}
Let $I_n=[a_n,b_n]$ be \textbf{closed, bounded} intervals that are
\textbf{nested}, i.e.\ $I_1\supseteq I_2\supseteq I_3\supseteq\cdots$.
Then
\begin{enumerate}[label=(\alph*),nosep]
  \item $\displaystyle\bigcap_{n=1}^\infty I_n\neq\emptyset$; in fact
        $\bigcap_n I_n=[\alpha,\beta]$ where $\alpha=\sup_n a_n$ and
        $\beta=\inf_n b_n$;
  \item if in addition $b_n-a_n\to0$, then $\bigcap_n I_n$ is a single
        point $\{\alpha\}$, and $a_n\to\alpha$, $b_n\to\alpha$.
\end{enumerate}
\end{thm}

\noindent\key{Picture.} The left endpoints move right, the right endpoints
move left, and they never cross. Completeness says there is no ``hole''
where they could meet, so some real number is trapped in all of the
intervals.

\begin{proof}[Full proof, step by step]
\emph{Step 1 (unpack ``nested'').} $I_{n+1}\subseteq I_n$ means
$[a_{n+1},b_{n+1}]\subseteq[a_n,b_n]$, i.e.
\[
a_n\le a_{n+1}\le b_{n+1}\le b_n\qquad\text{for all }n.
\]
So $(a_n)$ is increasing and $(b_n)$ is decreasing.

\emph{Step 2 (left endpoints never pass right endpoints).} Claim:
$a_m\le b_n$ for \emph{all} $m,n\in\N$ (not only $m=n$).
\begin{itemize}[nosep]
  \item If $m\le n$: $a_m\le a_n\le b_n$ ($a$ increasing, then
        $a_n\le b_n$).
  \item If $m>n$: $a_m\le b_m\le b_n$ ($a_m\le b_m$, then $b$
        decreasing).
\end{itemize}

\emph{Step 3 (define $\alpha$ by completeness).} Let
$A=\{a_m:m\in\N\}$. It is nonempty, and by Step 2 (with $n=1$) $b_1$ is an
upper bound. By the Completeness Axiom, $\alpha=\sup A$ exists.

\emph{Step 4 ($\alpha\in I_n$ for every $n$).} Fix $n$.
\begin{itemize}[nosep]
  \item $\alpha$ is an upper bound of $A$, so $a_n\le\alpha$.
  \item By Step 2, $b_n$ is an upper bound of $A$. Since $\alpha$ is the
        \emph{least} upper bound, $\alpha\le b_n$.
\end{itemize}
So $a_n\le\alpha\le b_n$, i.e.\ $\alpha\in I_n$. Since $n$ was arbitrary,
$\alpha\in\bigcap_nI_n$, and the intersection is nonempty.

\emph{Step 5 (the whole intersection).} Similarly
$\beta=\inf\{b_n\}$ exists, and Step 4 shows $\alpha$ is a lower bound of
$\{b_n\}$, so $\alpha\le\beta$. If $x\in[\alpha,\beta]$, then
$a_n\le\alpha\le x\le\beta\le b_n$ for every $n$, so $x\in\bigcap I_n$.
Conversely, if $x\in\bigcap I_n$ then $a_n\le x\le b_n$ for all $n$, so
$x$ is an upper bound of $\{a_n\}$ (giving $\alpha\le x$) and a lower
bound of $\{b_n\}$ (giving $x\le\beta$). Hence $\bigcap I_n=[\alpha,\beta]$.

\emph{Step 6 (part (b): uniqueness).} Suppose $b_n-a_n\to0$ and
$x,y\in\bigcap I_n$. For every $n$, both $x$ and $y$ lie in
$[a_n,b_n]$, so $|x-y|\le b_n-a_n$. Let $\eps>0$. For large $n$,
$b_n-a_n<\eps$, so $|x-y|<\eps$. Since $\eps$ was arbitrary,
$|x-y|=0$, so $x=y$ (the ``$<\eps$ for all $\eps$'' trick). So
$\bigcap I_n=\{\alpha\}$.

\emph{Step 7 (endpoints converge).} $0\le\alpha-a_n\le b_n-a_n\to0$ and
$0\le b_n-\alpha\le b_n-a_n\to0$, so by the Squeeze Theorem
$a_n\to\alpha$ and $b_n\to\alpha$.
\end{proof}

\noindent\key{Short proof via MCT.} $(a_n)$ is increasing and bounded
above by $b_1$, so $a_n\to\alpha$ by MCT. For fixed $n$, $a_m\le b_n$ for
all $m$, and limits preserve $\le$, so $\alpha\le b_n$. Also
$a_n\le\alpha$ because an increasing sequence stays below its limit. So
$\alpha\in I_n$ for all $n$.

\begin{ex}[Every hypothesis is needed]\label{ex:nitcounter}
\leavevmode
\begin{itemize}[nosep]
  \item \textbf{Closed is needed:} $I_n=(0,\frac1n)$ are nested, bounded
        and open, but $\bigcap_n(0,\frac1n)=\emptyset$. If $x>0$ were in
        all of them, pick $n>\frac1x$ (Archimedean). Then $x\ge\frac1n$,
        so $x\notin I_n$. And $x\le0$ is in none of them.
  \item \textbf{Bounded is needed:} $I_n=[n,\infty)$ are nested and
        closed, but $\bigcap_n I_n=\emptyset$: any $x$ fails to be in
        $I_n$ once $n>x$ (Archimedean).
  \item \textbf{Nested is needed:} $[n,n+1]$ are closed and bounded but not
        nested, and no point is in all of them.
  \item \textbf{$\R$ (completeness) is needed:} in $\Q$, let $I_n$ be the
        rational points of $[q_n,\,q_n+10^{-(n-1)}]$, where $q_n$ is the
        decimal truncation of $\sqrt2$ to $n-1$ decimal places. Their
        intersection in $\R$ is $\{\sqrt2\}$, so in $\Q$ it is empty.
\end{itemize}
\end{ex}

\begin{ex}[Bisection proof of Bolzano--Weierstrass using NIT]\label{ex:bwbisect}
Let $|x_n|\le M$ for all $n$. We find a convergent subsequence without
using Lemma~\ref{lem:monosub}.

\begin{proof}
\emph{Step 1 (start).} Let $I_1=[-M,M]$. It contains $x_n$ for
\emph{infinitely many} $n$ (in fact all $n$).

\emph{Step 2 (halve and keep an infinite half).} Suppose $I_k$ has been
chosen, has length $\frac{2M}{2^{k-1}}$, and contains $x_n$ for infinitely
many $n$. Split it into its closed left and right halves. If each half
contained $x_n$ for only finitely many $n$, then $I_k$ would too. So at
least one half contains $x_n$ for infinitely many $n$. Call it $I_{k+1}$
(choose the left half if both work). Its length is $\frac{2M}{2^k}$, and
$I_{k+1}\subseteq I_k$.

\emph{Step 3 (pick indices).} Let $n_1=1$ (so $x_{n_1}\in I_1$). Given
$n_k$, the set $\{n:x_n\in I_{k+1}\}$ is infinite, so it contains some
$n>n_k$. Let $n_{k+1}$ be the smallest such $n$. Then
$n_1<n_2<\cdots$ and $x_{n_k}\in I_k$ for all $k$.

\emph{Step 4 (NIT).} The $I_k$ are closed, bounded and nested, and their
lengths $\frac{4M}{2^k}\to0$. By Theorem~\ref{thm:nit} there is a single
$L\in\bigcap_kI_k$.

\emph{Step 5 (convergence).} For each $k$, both $x_{n_k}$ and $L$ lie in
$I_k$, so $|x_{n_k}-L|\le\frac{4M}{2^k}\le\frac{4M}k$ (using $2^k\ge k$).
Given $\eps>0$, choose $K>\frac{4M}\eps$. For $k\ge K$,
$|x_{n_k}-L|<\eps$. So $x_{n_k}\to L$.
\end{proof}
\noindent\key{Compare.} The peak proof uses completeness through MCT. The
bisection proof uses it through NIT. MCT, NIT, B--W and the Cauchy
criterion are all ways of saying that \emph{$\R$ has no holes}.
\end{ex}

\begin{ex}[Using NIT to locate a point: $\sqrt2$ by bisection]
Let $I_1=[1,2]$. At each step keep the half $[a,b]$ with $a^2\le2\le b^2$:
$[1,1.5]$, $[1.25,1.5]$, $[1.375,1.5]$, $[1.375,1.4375],\dots$ The
intervals are nested and closed, with lengths $2^{-(n-1)}\to0$, so NIT
gives a unique $c$ in all of them. Since $a_n^2\le2\le b_n^2$ and
$a_n,b_n\to c$, the product rule gives $a_n^2,b_n^2\to c^2$. Limits
preserve $\le$, so $c^2\le2\le c^2$, i.e.\ $c^2=2$. This is another
proof that $\sqrt2$ exists in $\R$, and it is the same idea as the
bisection method for root-finding.
\end{ex}

% =====================================================================
\section{Common mistakes that cost points}
\begin{enumerate}
  \item \textbf{Writing $\eps$--$N$ proofs backwards.} Scratch work goes
        from the goal to $N$. The proof must go from ``Let $\eps>0$'' to
        ``$<\eps$''.
  \item \textbf{Choosing $N$ that depends on $n$.} $N$ may depend only on
        $\eps$ (and fixed constants).
  \item \textbf{Forgetting to show that the sup is an upper bound.} Both
        parts (i) and (ii) are needed.
  \item \textbf{Assuming the sup is in the set.} $\sup(0,1)=1\notin(0,1)$.
  \item \textbf{Using limit laws when a limit doesn't exist}, e.g.\
        splitting $\frac{\sin n}{n}$ as $\lim\sin n\cdot\lim\frac1n$.
  \item \textbf{Passing to the limit before proving convergence} in
        recursive sequences. MCT comes first, then solve $L=f(L)$.
  \item \textbf{Multiplying an inequality by something that might be
        negative.}
  \item \textbf{Induction without writing $P(n+1)$ first}, and then
        ``proving'' the wrong target.
  \item \textbf{Strict vs.\ non-strict:} limits preserve $\le$ but not $<$.
  \item \textbf{Cauchy estimate that depends on $m$.} The bound on
        $|x_m-x_n|$ must depend only on $n$ (or on $N$), e.g.\
        $\frac1n-\frac1m<\frac1n$. Drop the $m$-term before choosing $N$.
  \item \textbf{Checking only neighbors.} Showing $|x_{n+1}-x_n|\to0$ does
        \emph{not} prove Cauchy (Example~\ref{ex:sqrtn}).
  \item \textbf{Putting $L$ into a Cauchy proof.} If you are asked to show
        a sequence is Cauchy because you \emph{don't} know $L$, you may not
        use $L$.
  \item \textbf{NIT with open or unbounded intervals.} The theorem needs
        closed \emph{and} bounded intervals (Example~\ref{ex:nitcounter}).
\end{enumerate}

% =====================================================================
\newpage
\section{Practice problems}
\label{sec:practice}
Try these under exam conditions. Solutions are in the next section.
Problems~\ref{prob:c1}--\ref{prob:c8} are on Cauchy sequences and
Problems~\ref{prob:n1}--\ref{prob:n3} are on nested intervals.

\begin{prob}
Prove by induction: $1+3+5+\dots+(2n-1)=n^2$ for all $n\in\N$.
\end{prob}

\begin{prob}
Prove by induction: $2^n\ge n+1$ for all $n\in\N$.
\end{prob}

\begin{prob}
Prove that $\sqrt3$ is irrational.
\end{prob}

\begin{prob}
Prove $A\cap(B\cup C)=(A\cap B)\cup(A\cap C)$.
\end{prob}

\begin{prob}
Find $\sup$ and $\inf$ of $S=\left\{(-1)^n\left(1-\frac1n\right):n\in\N\right\}$,
say whether each is attained, and prove your claim about the sup.
\end{prob}

\begin{prob}
Find $\sup$ and $\inf$ of $T=\left\{\frac{n}{n+1}:n\in\N\right\}$ with
proof.
\end{prob}

\begin{prob}
Let $A$ be nonempty and bounded above, and $c>0$. Prove
$\sup(cA)=c\sup A$.
\end{prob}

\begin{prob}
Let $A\subseteq B$ be nonempty and bounded. Prove $\inf B\le\inf A\le\sup A\le\sup B$.
\end{prob}

\begin{prob}
Use the definition to prove $\displaystyle\lim_{n\to\infty}\frac{2n^2+1}{n^2+3}=2$.
\end{prob}

\begin{prob}
Use the definition to prove $\displaystyle\lim_{n\to\infty}\bigl(\sqrt{n+1}-\sqrt n\bigr)=0$.
\end{prob}

\begin{prob}
Prove $\displaystyle\lim_{n\to\infty}\frac{n^2}{n+1}=\infty$ from the
definition.
\end{prob}

\begin{prob}
Prove that if $x_n\to L$ then $|x_n|\to|L|$. Is the converse true?
\end{prob}

\begin{prob}
Prove that $x_n=\frac{(-1)^n n}{n+1}$ diverges.
\end{prob}

\begin{prob}
Prove: if $x_n\to L$ and $x_n\ge 0$ for all $n$, then $L\ge0$.
\end{prob}

\begin{prob}
Let $x_1=2$ and $x_{n+1}=\frac{x_n+3}{2}$. Prove that $(x_n)$ converges and
find its limit.
\end{prob}

\begin{prob}
Compute (limit laws allowed):
(a) $\lim\frac{5n^3-n}{2n^3+7}$ \quad
(b) $\lim\frac{n+1}{n^2+1}$ \quad
(c) $\lim\left(\frac23\right)^n$ \quad
(d) $\lim\sum_{k=0}^n\frac1{2^k}$ \quad
(e) $\lim\frac{\cos(n^2)}{\sqrt n}$.
\end{prob}

\begin{prob}[True or false: justify or give a counterexample]\
\begin{enumerate}[label=(\alph*),nosep]
  \item If $x_n>0$ for all $n$ and $x_n\to L$, then $L>0$.
  \item Every bounded sequence converges.
  \item Every convergent sequence is bounded.
  \item If $x_ny_n\to0$ then $x_n\to0$ or $y_n\to0$.
  \item If $(x_n)$ and $(y_n)$ both diverge, then $(x_n+y_n)$ diverges.
  \item Every nonempty subset of $\Q$ that is bounded above has a sup in $\Q$.
  \item If $\sup S$ exists and $\sup S\in S$, then $S$ has a maximum.
\end{enumerate}
\end{prob}

\begin{prob}
Let $A,B$ be nonempty and bounded above. Prove
$\sup(A\cup B)=\max\{\sup A,\sup B\}$.
\end{prob}

\begin{prob}
Find $\sup$ and $\inf$ of $\left\{\frac1n+\frac1m:n,m\in\N\right\}$.
\end{prob}

\begin{prob}
Prove that if $x_n\to0$ and $(y_n)$ is bounded, then $x_ny_n\to0$. Use it
to find $\lim\frac{(-1)^n n}{n^2+1}$.
\end{prob}

\begin{prob}
Prove that if $x_n\to\infty$ then $\frac1{x_n}\to0$ (from the definitions).
\end{prob}

\begin{prob}
Use the definition to prove $\displaystyle\lim_{n\to\infty}\frac{3n+1}{n-2}=3$
(consider $n\ge3$).
\end{prob}

\begin{prob}
Prove that $\frac{n}{2^n}\to0$.
\end{prob}

\begin{prob}
Solve $|2x-1|<|x+2|$.
\end{prob}

\begin{prob}
Prove: if $x_n\to L>0$, then $x_n>\frac L2$ for all sufficiently large $n$.
\end{prob}

\begin{prob}
Show that $S=\{x\in\R:x^2<5\}$ is bounded above, explain why $\sup S$
exists without knowing its value, and identify $\sup S$.
\end{prob}

\subsection*{Cauchy sequences}

\begin{prob}\label{prob:c1}
Prove directly from the definition that $x_n=\dfrac{3n+2}{n+1}$ is Cauchy.
\end{prob}

\begin{prob}\label{prob:c2}
Let $s_n=\sum_{k=0}^n\frac1{k!}$. Prove $(s_n)$ is Cauchy. (Hint: first show
$k!\ge2^{k-1}$ for $k\ge1$.)
\end{prob}

\begin{prob}\label{prob:c3}
Let $t_n=\sum_{k=1}^n\frac1{\sqrt k}$. Prove $(t_n)$ is not Cauchy.
\end{prob}

\begin{prob}\label{prob:c4}
Without looking at Lemma~\ref{lem:csub}: prove that if $(x_n)$ is Cauchy
and a subsequence $x_{n_k}\to L$, then $x_n\to L$.
\end{prob}

\begin{prob}\label{prob:c5}
Suppose $0<r<1$ and $|x_{n+1}-x_n|\le r^n$ for all $n$. Prove $(x_n)$ is
Cauchy (and hence converges).
\end{prob}

\begin{prob}\label{prob:c6}
Let $x_1=1$, $x_2=2$, and $x_{n+2}=\frac{x_n+x_{n+1}}2$. Prove $(x_n)$ is
Cauchy and find its limit. (Hint: look at $d_n=x_{n+1}-x_n$, and at
$x_{n+1}+\frac12x_n$.)
\end{prob}

\begin{prob}[True or false: justify or give a counterexample]\label{prob:c7}\
\begin{enumerate}[label=(\alph*),nosep]
  \item If $x_{n+1}-x_n\to0$, then $(x_n)$ is Cauchy.
  \item Every Cauchy sequence of rational numbers has a rational limit.
  \item Every subsequence of a Cauchy sequence is Cauchy.
  \item If $(x_n)$ is Cauchy, then $(|x_n|)$ is Cauchy. Is the converse
        true?
  \item Every bounded sequence is Cauchy.
  \item If $(x_n)$ is Cauchy and $a\le x_n\le b$ for all $n$, then
        $\lim x_n\in[a,b]$.
\end{enumerate}
\end{prob}

\begin{prob}\label{prob:c8}
Prove from the definition that every Cauchy sequence is bounded. Then give
a bounded sequence that is not Cauchy.
\end{prob}

\subsection*{Nested intervals}

\begin{prob}\label{prob:n1}
Find $\bigcap_{n}\bigl[-\frac1n,\frac1n\bigr]$ and
$\bigcap_n\bigl(0,\frac1n\bigr]$, with proof. Why doesn't the second one
contradict the Nested Interval Theorem?
\end{prob}

\begin{prob}\label{prob:n2}
Let $I_n=\bigl[1-\frac1n,\ 1+\frac1{n^2}\bigr]$. Show the intervals are
nested and find $\bigcap_nI_n$.
\end{prob}

\begin{prob}\label{prob:n3}
Let $I_n=[a_n,b_n]$ be closed, nested, with $b_n-a_n\to0$, and let $\alpha$
be the point in $\bigcap_nI_n$. If $x_n\in I_n$ for every $n$, prove
$x_n\to\alpha$.
\end{prob}

% =====================================================================
\newpage
\section{Solutions}
\label{sec:solutions}

\paragraph{1.} Base case: $1=1^2$. Target $P(n+1)$:
$1+3+\dots+(2n-1)+(2n+1)=(n+1)^2$. By the hypothesis the left side is
$n^2+2n+1=(n+1)^2$. \qed

\paragraph{2.} Base case: $2\ge2$. Assume $2^n\ge n+1$. Then
$2^{n+1}=2\cdot2^n\ge2(n+1)=2n+2\ge n+2$. \qed

\paragraph{3.} Suppose $\sqrt3=p/q$ in lowest terms, so $p^2=3q^2$. If
$3\nmid p$, then $p=3k\pm1$ and $p^2=9k^2\pm6k+1=3(3k^2\pm2k)+1$ is not a
multiple of 3, a contradiction. So $p=3k$, which gives $9k^2=3q^2$,
$q^2=3k^2$, and so $3\mid q$. This contradicts lowest terms. \qed

\paragraph{4.} ($\subseteq$) Let $x\in A\cap(B\cup C)$. Then $x\in A$, and
$x\in B$ or $x\in C$. Case $x\in B$: $x\in A\cap B$. Case $x\in C$:
$x\in A\cap C$. Either way $x$ is in the union.
($\supseteq$) If $x\in A\cap B$, then $x\in A$ and $x\in B\subseteq B\cup C$.
The case $x\in A\cap C$ is the same. \qed

\paragraph{5.} The elements are $0,\frac12,-\frac23,\frac34,-\frac45,\dots$.
Even $n$ gives $1-\frac1n\in[\frac12,1)$. Odd $n$ gives
$-(1-\frac1n)\in(-1,0]$. So $\sup S=1$ and $\inf S=-1$, and \emph{neither
is attained}.

\emph{Proof that $\sup S=1$:} every element satisfies
$(-1)^n(1-\frac1n)\le1-\frac1n<1$, so $1$ is an upper bound. Given
$\eps>0$, choose an even $n$ with $\frac1n<\eps$ (Archimedean; if $n$ is
odd, use $n+1$). Then $1-\frac1n\in S$ and $1-\frac1n>1-\eps$. \qed
(The inf is symmetric, using odd $n$.)

\paragraph{6.} $\frac{n}{n+1}=1-\frac1{n+1}$, which is increasing. So
$\inf T=\min T=\frac12$ (at $n=1$), and $\sup T=1$ is not attained.
Upper bound: $\frac n{n+1}<1$. Given $\eps>0$, pick $n$ with
$\frac1{n+1}<\eps$. Then $\frac{n}{n+1}>1-\eps$. \qed

\paragraph{7.} Let $\alpha=\sup A$. For $a\in A$: $a\le\alpha$, so
$ca\le c\alpha$ (as $c>0$), and $c\alpha$ is an upper bound of $cA$. Given
$\eps>0$, $\frac\eps c>0$, so there is $a\in A$ with $a>\alpha-\frac\eps c$.
Then $ca\in cA$ and $ca>c\alpha-\eps$. By Theorem~\ref{thm:epssup},
$\sup(cA)=c\alpha$. \qed

\paragraph{8.} Middle: pick any $a\in A$, so $\inf A\le a\le\sup A$.
Right: $\sup B$ is an upper bound of $B\supseteq A$, so it is an upper bound
of $A$. Hence $\sup A\le\sup B$ (sup is the \emph{least} upper bound).
Left: the same argument with lower bounds. \qed

\paragraph{9.}
$\left|\frac{2n^2+1}{n^2+3}-2\right|=\left|\frac{2n^2+1-2n^2-6}{n^2+3}\right|=\frac5{n^2+3}<\frac5{n^2}\le\frac5n$.
Let $\eps>0$ and choose $N>\frac5\eps$. For $n\ge N$:
$\left|\frac{2n^2+1}{n^2+3}-2\right|<\frac5n\le\frac5N<\eps$. \qed

\paragraph{10.}
$\sqrt{n+1}-\sqrt n=\frac{1}{\sqrt{n+1}+\sqrt n}<\frac1{\sqrt n}$. We want
$\frac1{\sqrt n}<\eps\iff n>\frac1{\eps^2}$. Let $\eps>0$ and choose
$N>\frac1{\eps^2}$. For $n\ge N$:
$|\sqrt{n+1}-\sqrt n-0|<\frac1{\sqrt n}\le\frac1{\sqrt N}<\eps$. \qed

\paragraph{11.} For $n\ge1$, $n+1\le2n$, so $\frac{n^2}{n+1}\ge\frac{n^2}{2n}=\frac n2$.
Given $M$, choose $N>2M$. For $n\ge N$: $\frac{n^2}{n+1}\ge\frac n2\ge\frac N2>M$. \qed

\paragraph{12.} By the reverse triangle inequality,
$\bigl||x_n|-|L|\bigr|\le|x_n-L|<\eps$ for $n\ge N$, with the same $N$ that
works for $x_n\to L$. The converse is \emph{false}: $x_n=(-1)^n$ has
$|x_n|\to1$ but diverges. (It is true when $L=0$, since $|x_n-0|=||x_n|-0|$.) \qed

\paragraph{13.} Suppose $x_n\to L$. The even terms are
$\frac{n}{n+1}\to1$ and the odd terms are $-\frac n{n+1}\to-1$. Directly:
take $\eps=\frac12$. For $n\ge N$ large, even $n$ gives $x_n>\frac12$ and
odd $n$ gives $x_n<-\frac12$, so these terms are more than $1$ apart. But
$|x_n-x_m|\le|x_n-L|+|L-x_m|<1$, a contradiction. (Or: two subsequences
with different limits.) \qed

\paragraph{14.} Suppose $L<0$. Take $\eps=-\frac L2>0$. For large $n$,
$x_n<L+\eps=\frac L2<0$, which contradicts $x_n\ge0$. \qed

\paragraph{15.} The terms are $x_1=2$, $x_2=2.5$, $x_3=2.75$, so the guess
is increasing with limit $3$.
\emph{Bound $x_n<3$:} true for $n=1$; if $x_n<3$ then $x_{n+1}<\frac{3+3}2=3$.
\emph{Increasing:} $x_{n+1}-x_n=\frac{3-x_n}{2}>0$ by the bound.
By MCT, $x_n\to L$. Taking limits in $x_{n+1}=\frac{x_n+3}2$ gives
$L=\frac{L+3}2$, so $L=3$. \qed

\paragraph{16.} (a) $\frac52$ \quad (b) $0$ (lower degree on top) \quad
(c) $0$ since $|\frac23|<1$ \quad (d) $\frac1{1-1/2}=2$ \quad
(e) $0$, since $\left|\frac{\cos(n^2)}{\sqrt n}\right|\le\frac1{\sqrt n}\to0$ (squeeze).

\paragraph{17.}
\begin{enumerate}[label=(\alph*),nosep]
  \item False: $\frac1n>0$ but $\frac1n\to0$. (You only get $L\ge0$.)
  \item False: $(-1)^n$. (It \emph{does} have a convergent subsequence, by B--W.)
  \item True: the ``convergent $\Rightarrow$ bounded'' theorem.
  \item False: $x_n=1,0,1,0,\dots$ and $y_n=0,1,0,1,\dots$ have $x_ny_n=0$,
        but neither converges to $0$.
  \item False: $x_n=(-1)^n$ and $y_n=-(-1)^n$ sum to $0$.
  \item False: $\{q\in\Q:q^2<2\}$ has no sup in $\Q$, since the sup would be
        $\sqrt2$. This is exactly why $\R$ needs the completeness axiom.
  \item True: $\sup S$ is an upper bound in $S$, which is the definition of
        a maximum.
\end{enumerate}

\paragraph{18.} Let $\alpha=\sup A$, $\beta=\sup B$, $M=\max\{\alpha,\beta\}$.
\emph{Upper bound:} if $x\in A\cup B$, then $x\le\alpha\le M$ or
$x\le\beta\le M$. \emph{Least:} say $M=\alpha$. Given $\eps>0$, there is
$a\in A\subseteq A\cup B$ with $a>\alpha-\eps=M-\eps$. \qed

\paragraph{19.} $\sup=\max=2$ (at $n=m=1$), since $\frac1n+\frac1m\le1+1$.
$\inf=0$, not attained: $0$ is a lower bound since every element is $>0$.
Given $\eps>0$, pick $n$ with $\frac1n<\frac\eps2$. Then
$\frac1n+\frac1n<\eps$, with $m=n$. \qed

\paragraph{20.} Let $|y_n|\le M$ with $M>0$. Given $\eps>0$, choose $N$ with
$|x_n|<\frac\eps M$ for $n\ge N$. Then $|x_ny_n|\le M|x_n|<\eps$.
Application: $\frac{(-1)^nn}{n^2+1}=(-1)^n\cdot\frac{n}{n^2+1}$, where
$(-1)^n$ is bounded and $0<\frac n{n^2+1}<\frac1n\to0$. So the limit is
$0$. \qed

\paragraph{21.} Let $\eps>0$. Apply $x_n\to\infty$ with $M=\frac1\eps$:
there is $N$ with $x_n>\frac1\eps>0$ for $n\ge N$. Then
$0<\frac1{x_n}<\eps$, so $\left|\frac1{x_n}-0\right|<\eps$. \qed

\paragraph{22.} $\frac{3n+1}{n-2}-3=\frac{7}{n-2}>0$ for $n\ge3$. For
$n\ge3$ we have $n-2\ge\frac n3$ (equivalent to $\frac{2n}{3}\ge2$), so
$\frac7{n-2}\le\frac{21}n$. Given $\eps>0$, choose $N>\max\{3,\frac{21}\eps\}$.
For $n\ge N$: $\left|\frac{3n+1}{n-2}-3\right|\le\frac{21}n\le\frac{21}N<\eps$. \qed

\paragraph{23.} By the binomial theorem, $2^n\ge\binom n2=\frac{n(n-1)}2$ for
$n\ge2$, so $0<\frac n{2^n}\le\frac2{n-1}$. Given $\eps>0$, choose
$N>1+\frac2\eps$. For $n\ge N$: $\frac n{2^n}\le\frac2{n-1}<\eps$. \qed

\paragraph{24.} Both sides are $\ge0$, so the inequality is equivalent to
comparing squares:
$(2x-1)^2<(x+2)^2\iff3x^2-8x-3<0\iff(3x+1)(x-3)<0\iff-\frac13<x<3$.
(Alternatively, split into cases at $x=\frac12$ and $x=-2$.)

\paragraph{25.} Take $\eps=\frac L2>0$. There is $N$ with
$|x_n-L|<\frac L2$ for $n\ge N$, so
$x_n>L-\frac L2=\frac L2$. \key{This is the ``eventually positive'' lemma,
used in the quotient rule.} \qed

\paragraph{26.} If $x\ge3$ then $x^2\ge9>5$, so every $x\in S$ satisfies
$x<3$, and $3$ is an upper bound. $S\ne\emptyset$ since $0\in S$. By the
Completeness Axiom, $\sup S$ exists, even though we haven't computed it.
In fact $\sup S=\sqrt5$, whose existence was proved in lecture by exactly
this sup argument for $\sqrt2$.

\subsection*{Cauchy sequences}

\paragraph{\ref{prob:c1}.} $x_n=3-\frac1{n+1}$. For $m>n$:
$|x_m-x_n|=\frac1{n+1}-\frac1{m+1}<\frac1n$. Let $\eps>0$ and choose
$N>\frac1\eps$. For $m>n\ge N$: $|x_m-x_n|<\frac1n\le\frac1N<\eps$. \qed

\paragraph{\ref{prob:c2}.} \emph{Claim:} $k!\ge2^{k-1}$ for $k\ge1$.
Base: $1!=1=2^0$. Step: $(k+1)!=(k+1)k!\ge2\cdot2^{k-1}=2^k$. For
$m>n\ge1$:
\[
0<s_m-s_n=\sum_{k=n+1}^m\frac1{k!}\le\sum_{k=n+1}^m\frac1{2^{k-1}}
<\frac{1/2^n}{1-\frac12}=\frac1{2^{n-1}}\le\frac1n,
\]
using $2^{n-1}\ge n$ (Problem 2). Given $\eps>0$, choose $N>\frac1\eps$.
Then $|s_m-s_n|<\eps$ for $m>n\ge N$. \qed (The limit is $e$.)

\paragraph{\ref{prob:c3}.} Take $\eps=\frac12$. Given $N$, let $n=N$ and
$m=2N$. The sum $t_{2N}-t_N=\sum_{k=N+1}^{2N}\frac1{\sqrt k}$ has $N$
terms, each $\ge\frac1{\sqrt{2N}}$, so
$t_{2N}-t_N\ge\frac N{\sqrt{2N}}=\sqrt{\frac N2}\ge\frac1{\sqrt2}>\frac12$.
\qed

\paragraph{\ref{prob:c4}.} See Lemma~\ref{lem:csub}. The key line:
choose $N_1$ (Cauchy, $\frac\eps2$) and $K$ (subsequence, $\frac\eps2$),
fix $k=\max\{K,N_1\}$ so $n_k\ge k\ge N_1$. Then for $n\ge N_1$,
$|x_n-L|\le|x_n-x_{n_k}|+|x_{n_k}-L|<\eps$. \qed

\paragraph{\ref{prob:c5}.} For $m>n$, telescoping and the triangle
inequality give
$|x_m-x_n|\le\sum_{k=n}^{m-1}|x_{k+1}-x_k|\le\sum_{k=n}^{m-1}r^k<\frac{r^n}{1-r}$.
Since $r^n\to0$ (basic limit, $0<r<1$), given $\eps>0$ there is $N$ with
$r^N<\eps(1-r)$. For $m>n\ge N$, $r^n\le r^N$, so
$|x_m-x_n|<\frac{r^N}{1-r}<\eps$. \qed

\paragraph{\ref{prob:c6}.} Let $d_n=x_{n+1}-x_n$. Then
$d_{n+1}=\frac{x_n+x_{n+1}}2-x_{n+1}=-\frac12d_n$, and $d_1=1$, so
$|d_n|=\frac1{2^{n-1}}=2\bigl(\frac12\bigr)^n$. As in Problem~\ref{prob:c5}
(with an extra factor $2$),
$|x_m-x_n|<\frac{2(1/2)^n}{1/2}=\frac4{2^n}\le\frac4n$ for $m>n$, so
$(x_n)$ is Cauchy and converges to some $L$.

\emph{Finding $L$.} Taking limits in the recurrence gives $L=\frac{L+L}2$,
which is true for every $L$ and tells us nothing. Use the invariant instead:
$x_{n+2}+\frac12x_{n+1}=\frac{x_n+x_{n+1}}2+\frac12x_{n+1}=x_{n+1}+\frac12x_n$,
so $x_{n+1}+\frac12x_n=x_2+\frac12x_1=\frac52$ for all $n$. Letting
$n\to\infty$: $L+\frac L2=\frac52$, so $L=\frac53$. \qed

\paragraph{\ref{prob:c7}.}
\begin{enumerate}[label=(\alph*),nosep]
  \item False: $x_n=\sqrt n$ (Example~\ref{ex:sqrtn}).
  \item False: the decimal truncations of $\sqrt2$ are rational and
        Cauchy, but their limit is $\sqrt2\notin\Q$.
  \item True: if $|x_n-x_m|<\eps$ for $n,m\ge N$, then for $k,j\ge N$ we
        have $n_k\ge k\ge N$ and $n_j\ge N$, so $|x_{n_k}-x_{n_j}|<\eps$.
  \item True: $\bigl||x_n|-|x_m|\bigr|\le|x_n-x_m|$ (reverse triangle
        inequality), with the same $N$. The converse is false: $(-1)^n$.
  \item False: $(-1)^n$.
  \item True: it converges (Cauchy criterion), and limits preserve
        $a\le x_n\le b$.
\end{enumerate}

\paragraph{\ref{prob:c8}.} See Lemma~\ref{lem:cbdd}: use $\eps=1$ to get
$|x_n|<1+|x_N|$ for $n\ge N$, then
$M=\max\{|x_1|,\dots,|x_{N-1}|,1+|x_N|\}$. A bounded sequence that is
not Cauchy: $(-1)^n$. \qed

\subsection*{Nested intervals}

\paragraph{\ref{prob:n1}.} $\bigcap[-\frac1n,\frac1n]=\{0\}$: $0$ is in
every interval. If $x\ne0$, choose $n>\frac1{|x|}$; then $|x|>\frac1n$,
so $x\notin[-\frac1n,\frac1n]$. (Or apply NIT(b): the lengths $\frac2n\to0$.)

$\bigcap(0,\frac1n]=\emptyset$: if $x\le0$ it is in none of them. If
$x>0$, choose $n>\frac1x$, so $x>\frac1n$ and $x\notin(0,\frac1n]$.
This doesn't contradict NIT because the intervals are \emph{not closed}:
the point they shrink to, $0$, has been removed. \qed

\paragraph{\ref{prob:n2}.} \emph{Nested:} $1-\frac1n\le1-\frac1{n+1}$ and
$1+\frac1{(n+1)^2}\le1+\frac1{n^2}$, so $a_n\le a_{n+1}$ and
$b_{n+1}\le b_n$, which gives $I_{n+1}\subseteq I_n$. \emph{Intersection:} $1\in I_n$
for all $n$. The lengths $\frac1n+\frac1{n^2}\le\frac2n\to0$, so by
NIT(b) the intersection is exactly $\{1\}$. \qed

\paragraph{\ref{prob:n3}.} Both $x_n$ and $\alpha$ are in $[a_n,b_n]$, so
$|x_n-\alpha|\le b_n-a_n$. Given $\eps>0$, choose $N$ with
$b_n-a_n<\eps$ for $n\ge N$. Then $|x_n-\alpha|<\eps$ for $n\ge N$. \qed

% =====================================================================
\newpage
\section{Practice Exam 1 (75 minutes)}
\label{sec:mock}
Set a 75-minute timer, use no notes, and write full proofs. Solutions are
in the next section.

\begin{mprob}[Definitions, 10 pts]
State precisely: (a) $\sup S$; (b) $x_n\to L$; (c) $(x_n)$ is Cauchy, and
its negation; (d) the Completeness Axiom; (e) the Nested Interval Theorem.
\end{mprob}

\begin{mprob}[Induction, 15 pts]
Prove that $\displaystyle\sum_{k=1}^n\frac1{k^2}\le2-\frac1n$ for all
$n\in\N$.
\end{mprob}

\begin{mprob}[Sup/inf, 15 pts]\label{mock:recip}
Let $A\subseteq(0,\infty)$ be nonempty with $a=\inf A>0$. Let
$B=\{\frac1x:x\in A\}$. Prove $\sup B=\frac1a$.
\end{mprob}

\begin{mprob}[$\eps$--$N$, 15 pts]
Prove from the definition that
$\displaystyle\lim_{n\to\infty}\frac{n^2+2n}{3n^2-1}=\frac13$.
\end{mprob}

\begin{mprob}[Monotone sequences, 15 pts]
Let $x_1=1$ and $x_{n+1}=\sqrt{3+2x_n}$. Prove $(x_n)$ converges and find
its limit.
\end{mprob}

\begin{mprob}[Cauchy sequences, 20 pts]\label{mock:cauchy}
\leavevmode
\begin{enumerate}[label=(\alph*),nosep]
  \item Prove that every convergent sequence is Cauchy.
  \item Let $s_n=\sum_{k=1}^n\frac{\cos k}{k^2}$. Prove $(s_n)$ converges.
        (Note: $(s_n)$ is not monotone.)
  \item Let $u_n=\sum_{k=1}^n\frac1{2k-1}$. Prove $(u_n)$ is not Cauchy.
\end{enumerate}
\end{mprob}

\begin{mprob}[True/False, 10 pts]
Justify each answer briefly or give a counterexample.
\begin{enumerate}[label=(\alph*),nosep]
  \item If $x_n<1$ for all $n$ and $x_n\to L$, then $L<1$.
  \item If $(x_n)$ is increasing and not bounded above, then $x_n\to\infty$.
  \item If $(|x_n|)$ converges, then $(x_n)$ converges.
  \item $\sup\{r\in\Q:r<\sqrt2\}=\sqrt2$.
  \item If $(x_n)$ converges and $(y_n)$ diverges, then $(x_n+y_n)$
        diverges.
  \item If $I_n$ are nonempty, nested and bounded intervals, then
        $\bigcap_nI_n\ne\emptyset$.
\end{enumerate}
\end{mprob}

% =====================================================================
\newpage
\section{Practice Exam 1: solutions}

\paragraph{1.} (a)--(d): see the definitions cheat sheet. Negation of
Cauchy: $\exists\eps>0\ \forall N\ \exists n,m\ge N$ with
$|x_n-x_m|\ge\eps$. (e): see Theorem~\ref{thm:nit}. Say ``closed,
bounded, nested'' explicitly.

\paragraph{2.} $P(n)$: $\sum_{k=1}^n\frac1{k^2}\le2-\frac1n$. Base:
$1\le2-1$. Assume $P(n)$. Then
$\sum_{k=1}^{n+1}\frac1{k^2}\le2-\frac1n+\frac1{(n+1)^2}$, and we need
this to be $\le2-\frac1{n+1}$, i.e.
$\frac1{(n+1)^2}\le\frac1n-\frac1{n+1}=\frac1{n(n+1)}$. This holds because
$n(n+1)\le(n+1)^2$. \qed\
\key{Lesson:} trying to prove ``$\le2$'' directly by induction fails,
because the hypothesis is too weak. The stronger statement is easier to
prove.

\paragraph{3.} \emph{Upper bound:} for $x\in A$, $x\ge a>0$, so
$\frac1x\le\frac1a$. So $\frac1a$ is an upper bound of $B$ and
$\sup B\le\frac1a$.
\emph{Least} (Quiz 1 style, no $\eps$): let $u$ be any upper bound of $B$.
Then $u>0$, since $B$ has positive elements. For each $x\in A$,
$\frac1x\le u$, so $x\ge\frac1u$. So $\frac1u$ is a lower bound of $A$,
and $\frac1u\le a$ ($a$ is the \emph{greatest} lower bound). Hence
$u\ge\frac1a$. So $\frac1a$ is the least upper bound. \qed\
(If $\inf A=0$, then $B$ is unbounded.)

\paragraph{4.} $\frac{n^2+2n}{3n^2-1}-\frac13=\frac{6n+1}{3(3n^2-1)}>0$.
For $n\ge1$: $6n+1\le7n$ and $3n^2-1\ge2n^2$, so this is
$\le\frac{7n}{6n^2}=\frac7{6n}$. Let $\eps>0$ and choose
$N>\frac7{6\eps}$. For $n\ge N$:
$\bigl|\frac{n^2+2n}{3n^2-1}-\frac13\bigr|\le\frac7{6n}\le\frac7{6N}<\eps$.
\qed

\paragraph{5.} \emph{Bounds:} $0<x_n<3$ by induction. True for $n=1$,
and $0<x_n<3$ gives $3<3+2x_n<9$, so $0<\sqrt3<x_{n+1}<3$.
\emph{Increasing:} $x_{n+1}^2-x_n^2=3+2x_n-x_n^2=(3-x_n)(1+x_n)>0$, and
both terms are positive, so $x_{n+1}>x_n$.
By MCT, $x_n\to L$, and $L\ge x_1=1$. Square the recurrence:
$x_{n+1}^2=3+2x_n$. By limit laws, $L^2=3+2L$, so $(L-3)(L+1)=0$. Since
$L\ge1$, $L=3$. \qed

\paragraph{6.} (a) Part I of Theorem~\ref{thm:cauchy} (the $\frac\eps2$
argument).
(b) For $m>n$: $|s_m-s_n|=\bigl|\sum_{k=n+1}^m\frac{\cos k}{k^2}\bigr|
\le\sum_{k=n+1}^m\frac1{k^2}<\frac1n$, by
Example~\ref{ex:sumsq}. Choose $N>\frac1\eps$. So $(s_n)$ is Cauchy, and
by the Cauchy criterion it converges. (MCT doesn't apply, since the terms
$\cos k$ change sign.)
(c) Take $\eps=\frac14$. Given $N$, let $n=N$, $m=2N$. The sum
$u_{2N}-u_N=\sum_{k=N+1}^{2N}\frac1{2k-1}$ has $N$ terms, each
$\ge\frac1{4N-1}>\frac1{4N}$, so $u_{2N}-u_N>\frac14$. \qed

\paragraph{7.}
\begin{enumerate}[label=(\alph*),nosep]
  \item False: $1-\frac1n\to1$.
  \item True: given $M$, $M$ is not an upper bound, so some $x_N>M$. For
        $n\ge N$, $x_n\ge x_N>M$.
  \item False: $(-1)^n$.
  \item True: $\sqrt2$ is an upper bound. Given $\eps>0$, density gives
        $r\in\Q$ with $\sqrt2-\eps<r<\sqrt2$.
  \item True: if $x_n+y_n$ converged, then $y_n=(x_n+y_n)-x_n$ would
        converge.
  \item False: $(0,\frac1n)$ are not closed and have empty intersection.
        (True if the intervals are closed.)
\end{enumerate}

\end{document}
